% 本文件由 LaTeX 排版 Agent 根据有效 translation.json 自动生成。
% author=Novatian; work_id=0511; title=论圣三

\chapter{论圣三}

% structure: p0001 source=h1 target=omitted reason=page-title-already-rendered-by-parent

\EditorialNote{诺瓦替安有关圣三的论文分为三十一章。他首先从第一章到第八章，考虑信德或信德的规则中那些命令我们相信天主圣父和全能的主、万有绝对完美的创造者的话语。其中，除了其他神圣属性，他还部分地通过理性、部分地通过圣经将无限、永恒、合一、美善、不变、不死、精神性归于祂；并补充说，不能将情感或肢体归于天主，这些只是圣经中以拟人化的方式论及天主。}\par

% structure: p0003 source=h2 target=section reason=relative-heading-rank; base=h2

\section{第一章}

\Emph{诺瓦替安为论述圣三，首先从信德规则出发，阐明我们应首先相信天主圣父和全能的主，万物的绝对创始者。受造之工程被优美地描述。人的自由意志被肯定；天主在施加惩罚时显示的仁慈被展示；义人与不义之人死后灵魂的状态被确定。}\par

真理的规则要求我们首先相信天主圣父和全能的主；即万有的绝对完美的创造者，祂以崇高的威严悬垂诸天，以低沉的质体奠定大地，以流动的湿气散布海洋，并以合宜的器具装饰和供给这一切。因为在坚固的天穹中，祂唤起了带来光明的日出；祂在月盈中充满月亮的白色圆球作为夜晚的慰藉；此外，祂还用闪烁光芒的多样光辉点燃星芒；祂愿意这一切在合法的轨道上环绕世界的整个范围，从而为人类带来日、月、年、标志、季节以及其他种类的恩惠。在大地上，祂又将最高的山峰高举至顶峰，将山谷投下深渊，平整平原，为了人类的各种服务有益地规定了动物群。祂还为了人类未来使用的益处，设立了树林的橡树。祂发展出作为食物的收成。祂打开了泉源的出口，并将它们倾入流动的河流。在此之后，唯恐祂也不为眼睛的愉悦本身作准备，祂以花卉的各种颜色装饰了一切，以供观赏者之乐。此外，甚至在海中，尽管大海本身因其广度和效用而令人惊叹，祂也造了多样的受造物，有时体型适中，有时巨大，以祂设立的多样性见证创造者的智慧。并且，不满足于此，唯恐咆哮和急流的水以牺牲其人类拥有者为代价夺取陌生的元素，祂用海岸封闭了其界限；因此，当狂怒的巨浪和起泡沫的水从其深处而来时，它应回到自身，不越过其隐藏的边界，而遵守其规定法律，以便人更谨慎地遵守神圣法律，正如元素自身遵守他们一样。在此之后，祂又将人置于世界的首位，而人是按天主的肖像造的，祂赋予人理智、理性和预见，使人能模仿天主；虽然人身体的初始元素是属土的，但本质被天堂的和神圣的嘘气所激发。当祂赐予人一切为其服务时，祂愿意人独自是自由的。并且，唯恐无限制的自由陷入危险，祂订立了一条诫命，在其中人被教导，树的果实中没有恶；但人被预先警告，如果人因藐视所赐予的法律而行使自由意志，恶就会产生。因为，一方面，人必须自由，以免天主的肖像不适宜地处于束缚中；另一方面，必须加上法律，以免无节制的自由甚至爆发到对赐予者的藐视。因此，人可以根据其心智在任一方向的倾向，接受相称的赏报和应得的惩罚，在自己能力中有权选择做什么；由此，因嫉妒，死亡回到了他身上；看，虽然他可以通过服从逃避死亡，但他在邪恶的建议下急于成为神而冲入死亡。然而，天主仁慈地缓和了惩罚，诅咒的不是他的人，而是他在地上的劳作。而且，所要求的不在没有人的知识下发生；祂却彰显了人在基督里未来发现和救恩的希望。他之所以被禁止触摸生命树的果实，并非出于嫉妒的恶意毒害，而是唯恐在基督原先赦免其罪之前，他永远活着，将罪恶不朽的惩罚永远带着。此外，在更高的区域，即甚至在穹苍之上，那些我们眼睛如今无法看见的区域，祂预先规定了天使，安排了灵性权能，设立了座和权柄，并创建了诸天许多其他无限的空间和祂奥迹的无界工程；因此，这世界虽巨大，却几乎显得是最后而不是唯一的有形之物的工程。并且，确实，地下的东西本身也不是缺乏分布和安排的权能的。因为有一个地方，义人和不义之人的灵魂被带去，意识到来日审判的预期判决；因此我们能在一切方向看到天主工程的浩大，不限于这世界的怀抱（尽管正如我们所说，这怀抱是广阔的），而且也能在世界本身的深渊和深处之下构想。如此，考虑到工程的浩大，我们应恰当地赞美如此结构的创造者。\par

% structure: p0006 source=h2 target=section reason=relative-heading-rank; base=h2

\section{第二章}

\Emph{天主超越万有，自身包含万有，无限、永恒，超越人的心智；言语无法解释，高于一切崇高。}\par

在这一切之上，祂本身，包含万有，在祂之外没有任何虚空，没有为更高的天主（如某些人所想象的那样）留下余地。因为，事实上，祂已将万有都包含在完美伟大和能力的怀抱中，祂始终专注于自己的工作，渗透万有，推动万有，赋予万有生命，鉴察万有，并将不协调的材料联结成所有元素的和谐，以致于从这些不同的原则中，通过协和的统一，一个世界得以如此建立，以至于除非唯独创造它的祂命令它解散（为了赐给我们别的、更伟大的事物），否则没有任何力量可以解散它。因为我们读到祂包含万有，因此，在祂之外不可能有任何东西。因为，既然祂没有任何开端，所以自然也意识不到终结；除非或许——但愿不会有这种想法——祂在某个时间开始存在，并且不在万有之上，而是祂在别的事物之后开始存在，那么祂就低于在祂之前的存在，从而被发现能力较小，因为即使就时间本身而言，祂也被认为是后来的。因此，祂始终是无边的，因为没有比祂更伟大的；始终是永恒的，因为没有比祂更古老的。因为没有开端的东西不能被任何事物所先于，因为祂没有时间。因此祂是不朽的，因为祂的完整性不会因任何终结而结束。既然没有开端的东西就没有规律，祂排除了时间模式，因为祂感觉不欠任何人的债。因此，关于祂，以及关于那些出于祂且在祂之内的事物，人的心智无法恰当地构想它们是什么、多么伟大、像什么；人的言语雄辩也无法展现出接近祂威严水平的力量。因为要构想和述说祂的威严，正如所有雄辩理所当然地沉默，所有心智都贫乏。因为祂比心智本身更伟大；无法构想祂有多么伟大，因为如果能被构想，祂就会比能够构想祂的人心更小。此外，祂比所有言语都更伟大，无法被宣告；因为如果能被宣告，祂就会比人的言语更小，通过宣告，祂可以被包围和包含。因为任何能关于祂被想到的都必须小于祂本身；任何能被宣告的，当与祂本身比较时，都必须小于祂。此外，我们可以在某种程度上在沉默中意识到祂，但我们无法在话语中像祂那样展开祂。因为如果你称祂为\Emph{光}，你是在谈论祂的受造物而非祂本身——你并没有宣告祂；或者如果你称祂为\Emph{力量}，你更是在谈论和展现祂的能力而非谈论祂本身；或者如果你称祂为\Emph{威严}，你更是在描述祂的尊荣而非祂本身。我为什么要长篇累牍地逐一列举祂的属性呢？我将立即展开全部。无论你在任何方面可能宣告祂的任何东西，你更是在展开祂的某种情况和能力而非祂本身。因为对于那比所有言语和思想都更伟大的那一位，你能合适地说或想什么呢？除了一种方式——我们如何能做到？我们怎么可能构想我们如何把握这些事本身？——如果我们思考祂是那种无法在性质或数量上被理解、甚至无法进入思想本身的东西，我们就会在心灵上把握天主是什么。因为如果我们眼睛的敏锐在看太阳时变得迟钝，以至于视线被迎面而来的光线亮度所压倒，无法看见太阳本身，那么我们心智的敏锐在我们所有关于天主的思考中也会遭受同样的情况，而且我们越直接地努力思考天主，心智本身就越多地被它自己思想的光所蒙蔽。因为——再重复一次——对于那高于所有崇高、高于所有高度、深于所有深度、清于所有光明、亮于所有明亮、灿烂于所有光辉、强于所有力量、有力于所有能力、强壮于所有力量、大于所有威严、有能于所有潜能、富于所有财富、智慧于所有智慧、仁慈于所有慈爱、善良于所有良善、正义于所有正义、慈悲于所有怜悯的那一位，你能恰当地说什么呢？因为所有种类的美德必然小于祂本身，祂既是所有美德的天主和父，因此可以真实地说，天主是那种没有任何事物能与祂相比的存有。因为祂高于所有可言说之物。因为祂是一种心智，生成并充满万有，没有任何时间开端或终结，以最高和最完美的理性，控制着事物自然联结的原因，以造福众生。\par

% structure: p0009 source=h2 target=section reason=relative-heading-rank; base=h2

\section{第三章}

\Emph{天主是万物的创造者、他们的主与父，这由圣经证明。}\par

我们承认并知道祂即是天主，万物的创造者——因祂的大能而是主，因祂的训导而是父——我说的是祂，那发言而万物造成，发令而万物生出的祂。关于祂记载说：\BibleQuote{祢以智慧造了万物；} 梅瑟论及祂说：\BibleQuote{天主在天上，在地下一并；} \BibleRef{申 4:39} 祂按依撒意亚说：\BibleQuote{祂以手掌量过天空，以手心量过大地；} \BibleQuote{祂垂视大地，使大地震动；祂划定大地的圆圈，地上的居民好像蚱蜢；祂用天平称过山岳，用秤称过丛林，} 这就是说，凭着神圣安排的准确考验；如果不用相等的重量来平衡，就容易崩塌；祂以公正平衡了这尘世重担。祂藉先知说：\BibleQuote{我是天主，没有别的神} \BibleRef{依 45:22} 祂藉同一先知说：\BibleQuote{因为我不将我的荣耀给予别的神，} \BibleRef{依 13:8} 这是为了排除所有外邦人和异端者及他们的虚构；证明手工艺者手造的不是神，异端者心智虚构的也不是神。因为须询问匠人才存在的，不是神。祂又藉先知补充说：\BibleQuote{上天是我的宝座，下地是我的脚凳：你们要为我建筑什么殿宇，哪里是我安息的地方？} 这是要表明，那世界不能容纳的，更不会被殿宇容纳；祂说这些不是为了自夸，而是为了我们的知识。因为祂不希望我们从祂的伟大中得到荣耀，而是愿意像父亲一样，赐予我们虔诚的智慧。祂还希望吸引我们粗野、傲慢、因粗鲁的残暴而顽固的心灵趋向温柔，便说：\BibleQuote{我的圣神将安息在谁身上？除非是在那谦卑、安静、对我言语战栗的人身上。} \BibleRef{依 66:2} ——这样，人在某种程度上可以认识天主是多么伟大，因所赐予的圣神而学习敬畏祂。祂同样愿意更深地进入我们的知识，并通过激励我们的心神转向敬拜祂，说：\BibleQuote{我是上主，造了光，也创造了黑暗。} 这样我们就不会认为某种本性——我不知道是什么——是那控制昼夜交替的造物主，而是更真实地承认天主是它们的创造者。既然无法以肉眼看见祂，我们就当从祂工程的伟大、能力和威严中认识祂。保禄宗徒说：\BibleQuote{祂那看不见的事，} 自从创世以来，藉着所造之物就明明可见，即祂永远的大能和神性；这样人的心智，从明显的事学习隐藏的事，从所应观看的工程之伟大，能以心眼思量那建造者的伟大。同一宗徒论及祂说：\BibleQuote{愿尊崇和荣耀归于永世的君王，那不朽、不可见、独一的天主。} \BibleRef{弟前 1:17} 经上说：\BibleQuote{因为万有都是出于祂、藉着祂并归于祂。} \BibleRef{罗 11:33} 因为万有都是因祂的命令而存在，因为它们是\Emph{出于祂}的；并因祂的圣言而秩序井然，因为是\Emph{藉着祂}的；万有都回到祂的审判，因为在祂内期望自由，当朽坏被除去时，它们似乎被召回\Emph{归于祂}。\par

% structure: p0012 source=h2 target=section reason=relative-heading-rank; base=h2

\section{第四章}

\Emph{此外，祂是善的，永远相同，不变不化，独一无限；祂的名字无法宣告，祂不朽不灭。}\par

主惟独正确宣称那善的，其善行由整个世界作证；祂若不善，就不会创造这个世界。因为如果\Quote{样样都很好}\BibleRef{创 1:31}，那么，合理而言，受造之物既证明创造者是善的，而善的创造者之作品也必然是善的；因此，一切邪恶都是背离天主。因为祂自称是\Quote{完美的}、既是父又是审判者，尤其是一切恶行的惩罚者与审判者，祂绝不可能成为任何恶行的发起者或设计者。此外，邪恶临于人，并非出于其他原因，而是由于人背离了善的天主。而且，这一点在人身上是特定的，并非出于必然，而是人自己愿意如此。由此，什么是邪恶也显明了；并且，为了不让人以为天主有嫉妒，邪恶的来源也清楚了。因此，祂始终如一，从不改变或转化为任何形式，以免因变化而显得有死。因为从一物转向另一物的变化，被视为某种死亡的一部分。因此，祂从未有任何部分或尊荣的增加或增长，以免祂的圆满似乎曾有所欠缺；祂也未曾遭受任何损失，以免祂似乎承受了某种程度的死亡。祂过去是什么，永远是；祂是谁，永远是祂自己；祂具有什么品性，永远具有。因为增长意味着开端，损失则证明死亡与朽坏。因此祂说：\Quote{我是天主，我不改变}\BibleRef{拉 3:6}，因为那非受生的不能遭受变化，祂常保持其状态。因为凡在祂内构成天主性的一切，必然永远存在，凭自身的能力维持自己，使祂永远是天主。因此祂说：\Quote{我是自有者。}因为祂之所是具有这名号，因为它永远保持其自身的同一品质。因为变化会取消\Quote{我是自有者}这名号的力量；因为任何事物一旦变化，就在变化的那个具体方面显示出有死性：它不再是它过去所是的，因而开始成为它不是的；因此，合理而言，在天主内，祂的地位永远不变，因为祂没有因变化而遭受任何损失，始终与自身相似且相等。那非受生的不能改变：因为只有受造或受生之物才会变化；那些一度不存在之物，借着出生而学会存在，因此因受生而遭受变化。而那些既无受生也无创造者的事物，已将变化的能力排除在外，因为它们没有变化之因的开端。因此，祂被宣告为独一，没有同等者。因为凡能成为天主的，作为天主，必然是至高者。但凡是至高者，其至高必然达到没有任何同等者的程度。因此，那无可加添、无与伦比者必然是唯一且独一的；因为不可能有两个无限者，这是事物的本性所决定的。那无限者既无始也无终。因为那充满万有的，排除了另一个的开端。如果祂不包含一切所有——因为那被包含于某物之中的，必定小于那包含它的——祂就不再是天主，而会降服于另一个的权能之下，被那更大的所包含。因此，那个包含祂的，反而更应被称为天主。由此，天主的名号也无法被宣告，因为祂无法被理解。能被名称所包含的，是从其本性的状态中，以某种方式被理解的东西。因为名称是对那能被名称理解之物的指称。但当所讨论的对象如此超越，甚至连理智本身也不能恰当地将其收摄于一个形式时，又岂能通过一个称谓的单一词语恰当地表达出来呢？因为既然它超越理智，也必然超越称谓的意指能力。因此，当祂出于某些理由和场合，应用并偏爱祂的天主之名时，我们知道，这所展示的与其说是称谓的合法专有性，不如说是为它确定的某种意指性，人们投奔这名称，似乎就能借此获得天主的仁慈。因此，祂是不死不朽、不可朽坏的，不觉任何损失，也无终结。因为祂是不可朽坏的，所以是不死不朽的；因为祂是不死不朽的，所以当然也是不可朽坏的——二者彼此交替相涉，自身与自身相互联结，并因交替的连锁而延伸至永恒的状态：不死不灭来自不可朽坏，而不可朽坏也来自不死不灭。\par

% structure: p0015 source=h2 target=section reason=relative-heading-rank; base=h2

\section{第五章}

\Emph{如果我们探究圣经所描述的天主的忿怒、愤慨与憎恨，必须记住，这些不可被理解为具有人类恶习的特性。}\par

此外，如果我们读到祂的愤怒，并思考某些关于祂义怒的描述，且得知憎恨也被归于祂，我们却不可理解为它们是在人身上那样的恶习的含义上被归于祂。因为所有这些事物，尽管可能败坏人类，却丝毫不能败坏天主的能力。因为像这样的情绪，说它们在人内是正确的，而说它们在天主内则是不正确的。因为人可能被这些事物败坏，因为他会败坏；天主不会被它们败坏，因为祂不能败坏。这些事物固然有它们所能施加的力量，但仅限于在有能受影响之质料在先的地方，而不是在不能受影响之实体在先的地方。因为天主发怒，并非出于祂内的任何恶习。而是为我们的益处祂才如此；因为祂甚至在威胁时也是仁慈的，因为借着这些威胁，人们被唤回正义。因为恐惧对于那些缺乏向善动机的人是必要的，好叫那些舍弃理性的人至少能被恐惧所动。因此，所有这些，无论是天主的愤怒还是憎恨，或任何此类事物，为我们的医治而显示出来——正如情形所示——乃是出于智慧，而非出于恶习，也并不源于软弱；因此它们也不能导致天主的败坏。因为在我们内，构成我们身体的材质多样性，通常会激起那败坏我们的愤怒的不和；但这种不和，无论是本性的还是缺陷的，都不能存在于天主内，因为祂确实绝非由任何身体部分的结合所构成。因为祂是单纯的，没有任何物质混合，完全是那本质——不论它是什么，惟有祂知道——构成祂的存在，因为祂被称为神。因此，那些在人内是错误和败坏的事物，既然源于身体和质料本身的易朽性，在天主内就不能施加易朽的力量，因为，如我们所说，它们不是出于恶习，而是出于理性。\par

% structure: p0018 source=h2 target=section reason=relative-heading-rank; base=h2

\section{第六章}

\Emph{而且，虽然圣经常将天主的面貌改变为人的形像，但天主威严的度量却不包含在我们的身体本性的这些轮廓之内。}\par

虽然天上的圣经常常将天主的形象转变为人的样式——如它说：\Quote{主的眼睛注视义人}；或如它说：\Quote{上主天主闻到了馨香的气味}\BibleRef{创 8:21}；或如梅瑟领受的约版是\Quote{用天主的手指写的}\BibleRef{出 31:18}；或如以色列子民离开埃及地是\Quote{用大能的手和伸展的臂膀}；或如它说：\Quote{上主的口说了这些话}\BibleRef{依 1:21}；或如大地被立为\Quote{天主的脚凳}；或如它说：\Quote{侧耳倾听}\BibleRef{编下 19:16}——我们这些说法律是属神的人，并不认为这些我们身体本性的轮廓之中包含任何天主威严的方式或形态，而是将那种无限伟大的特征（可谓如此）散布于其平原之上，没有任何界限。因为经上记载：\Quote{我若升到天上，祢在那里；我若下到地狱，祢也在那里；我若展翅飞到海极，祢的手必拉住我，祢的右手必握住我。}因为我们根据圣经安排的尺度，承认其计划。因为那时先知还在按照信德的时期，用比喻论及天主，不是按天主本身，而是按百姓能接受的方式。因此，这类关于天主的话，不应归咎于天主，而应归咎于百姓。因此，百姓被允许建造会幕，然而天主并不被局限在会幕的围墙之内。同样，圣殿被建造，然而天主绝不受圣殿的限制。所以受到限制的不是天主，而是百姓的认知；并不是天主受束缚，而是百姓的理解力被视为有限。最后，主在福音中说：\Quote{时候将到，你们不在这山上，也不在耶路撒冷朝拜父}\BibleRef{若 4:21}，并说明了原因：\Quote{天主是神，因此朝拜的人必须用心灵和真理朝拜}\BibleRef{若 4:24}。因此，神圣的行动通过肢体展现；描述的并不是天主的形象或身体轮廓。因为当说起眼睛时，意指祂看见一切；当说起耳朵时，表示祂听到一切；当说起手指时，启示了祂意志的某种能力；当说起鼻孔时，显示祂认出祈祷如馨香之气；当说起手时，证明祂是万有的创造者；当说起手臂时，宣告没有任何本性能够抵抗祂手臂的能力；当说起脚时，揭示了祂充满万有，没有任何地方没有天主。因为祂不需要肢体或肢体的功能，因祂独一的判断，即使不出声，万有也侍奉并呈现在祂面前。祂为何需要眼睛？祂本身就是光。祂为何需要脚？祂无所不在。祂为何想要移动？没有任何地方在祂之外可以去。祂为何需要手？祂的意志，即使沉默，也是万有根基的建筑师。祂不需要耳朵，祂知道甚至未表达的意愿；祂需要舌头作什么？祂的念头就是命令。这些肢体对人而言是必需的，对天主则不然，因为人的计划若没有身体实现思想就会落空。此外，它们对天主不是必需的，因为祂的意志成就工作，不费任何气力，甚至工作与意志同时进行。而且，祂自己全是眼睛，因为祂全看见；全是耳朵，因为祂全听见；全是手，因为祂全工作；全是脚，因为祂全在万有之中。因为无论祂是什么，祂都是同一的。祂全是相等的，全在万有之中。因为祂在自己内没有多样性，而是单纯的。因为那些源于出生且走向消亡的东西，才被归结为肢体的多样性。而非质合之物不能意识到这些。那不朽的，无论它是什么，它本身是一、单纯且永恒的。因此，因为它是一，就不能被分解；因为凡是置于分解要求之外的东西，就摆脱了死亡的法则。\par

% structure: p0021 source=h2 target=section reason=relative-heading-rank; base=h2

\section{第七章}

\Emph{此外，当天主被称为神、光明、光时，天主并未被这些称呼充分表达。}\par

但是，当主说天主是神时，我认为基督如此论及圣父，是希望人理解比天主是神更深的含义。因为，虽然在祂的福音中，祂为增进人的理解力而推理，但祂自己仍以人尚能听见和接受的方式向人谈论天主；尽管如我们所说，祂如今正努力为听众增添对天主的宗教知识。因为我们发现经上记载：天主被称为\Emph{爱}，但由此并未宣告天主的实体是\Emph{爱}；祂被称为\Emph{光}，但其中并非天主的实体。然而，凡如此论及天主的，都是尽可能而言的，因此，当祂被称为神时，也并非以祂的全部所是而如此称呼；而是使人藉着理解在精神上日益进步，甚至达到神本身，且已在精神上改变，从而藉着神猜测天主是更为伟大的存在。因为那按其本然的存在，既不能以人的言语表达，也不能被人的耳朵接收，更不能被人的感知领会。因为\Quote{天主为爱祂的人所预备的，是眼睛未曾看见，耳朵未曾听见，人心也未曾想到的；}\BibleRef{格前 2:9}那么，许诺这一切的祂本身是何等伟大，人的心智和本性在领悟时都已失败！最后，如果你将神当作天主的实体，你便会使天主成为受造物。因为每个神都是受造物。因此，天主将变成被造之物。同样，如果依据梅瑟，你将天主接受为\Emph{火}，并说祂是受造物，你便是在宣告那被造之物，而非教导造物主。但这些更多是作为比喻使用，而非事实。因为，在旧约中，天主之所以被称为\Emph{火}，是为了使有罪的百姓心生恐惧，向他们提示一位审判者；同样，在新约中，祂被宣告为神，以便作为那些死于罪恶者的更新者和创造者，藉此仁慈的恩宠为信者所证实。\par

% structure: p0024 source=h2 target=section reason=relative-heading-rank; base=h2

\section{第八章}

\Emph{因此，这就是教会所认识和朝拜的天主；祂所眷顾和治理的万有，无论有形无形，都时时处处为祂作证。}\par

所以，教会知道并钦崇这位天主，弃绝了异端者的神话和虚构；宇宙和整个本性，无论可见的与不可见的，都为此作证；天使朝拜祂，星辰惊叹祂，海洋祝福祂，大地崇敬祂，地底下的一切都仰望祂；人的整个心智都意识到祂，即使它不表达\Emph{自身}；万物都因祂的命令而运动，泉水涌出，河流奔流，波浪兴起，一切受造物生育后代，风被约束吹拂，降雨降下，海洋被搅动，各处的一切都散发它们的丰饶。祂为永生的原祖特别安排了一个美丽的东方乐园；祂栽种了生命树，并在其旁边同样放置了另一棵知善恶树，赐下命令，并判决了罪的惩罚；祂使最正义的\Person{诺厄}因他的无罪和信德的功劳免于洪水的危险；祂提升了\Person{哈诺客}；祂拣选了\Person{亚巴郎}进入祂友谊的团体；祂保护了\Person{依撒格}；祂增加了\Person{雅各伯}；祂赐予\Person{梅瑟}为人民的领袖；祂拯救了叹息的以色列子民脱离奴役的轭；祂书写了法律；祂将我们祖先的后裔带入应许之地；祂藉祂的圣神教导了先知，并藉他们全体预许了祂的圣子基督；并在祂曾约定赐下祂的时刻，祂派遣了祂，并愿意藉着祂使我们认识祂，并将祂丰富的恩宠倾注在我们身上，将祂丰沛的圣神赐给贫穷卑微的人。并且，因为祂出于自己的自由意志既慷慨又仁慈，唯恐这整个大地因转离祂恩宠的溪流而枯萎，祂愿意使宗徒作为我们家族的奠基者，由祂的圣子派遣到全世界，以使人类的情况认识到它的创造者；并且，如果它选择跟随祂，就可能有一位它甚至在祈祷中如今可以称祂为圣父而非天主的。祂的眷顾在人类中，不仅对个人，而且对城市和国家本身，其毁灭已由先知的话语宣告，已经并正在进行；甚至遍及整个世界本身；祂已注定其结局、其苦难、其荒废和因不信而遭受的痛苦。此外，免得有人以为天主如此不懈的眷顾甚至不达到最微小的事物，主说：\Quote{两只麻雀中的一只，}主说，\Quote{不会没有圣父的旨意而掉落；但就是你们的头发也都被数过了。}\BibleRef{玛 10:29}祂的关怀和眷顾甚至不容以色列子民的衣服穿破，也不容他们脚上最破旧的鞋子磨损；此外，甚至被掳的青年人的外衣也不被烧毁。这并非没有理由；因为如果祂拥抱万物，包含万物——而万物和整体由个体组成——祂的关怀因而也将延伸到每一个个体，因为祂的眷顾达到整体，无论它是什么。因此，祂也坐在革鲁宾之上；也就是说，祂主宰祂工程的各种变种，那些控制其余受造物的生物被置于祂的宝座之下：一个水晶的覆盖物覆盖了一切；也就是说，天覆盖了一切，这天空曾按天主的命令从水的流动物质中凝结成穹苍，以便那坚固的硬度分开以前覆盖大地的水的中间，像背负一样支撑着其上水的重量，其力量由寒霜建立。而且，下面还有轮子——也就是说，季节——世界的所有成员藉此不断滚动；加上这些脚，使那些事物不永远静止，而是前进。而且，它们的四肢上布满了眼睛；因为天主的工程必须以永远警醒的检查来默观：在这些事物的中心，有炭火在其中，或者是因为我们这个尘世正趋向于审判的烈火之日；或者是因为天主的全部工程都是火热的，不是幽暗的，而是繁荣的。再者，唯恐因为这些事物源于地上的开始，它们由于起源的僵硬而自然不活动，内部精神的热本性被加到万物中；并且这种与寒冷物体结合的本性，为生命之目的，为一切提供均等的度量。因此，按\Person{达味}所说，这是天主的战车。他说：\Quote{天主的战车，}\Quote{增添了一万倍；}也就是说，是不可数、无限、无量的。因为，在赋予万物的自然法律的轭下，有些事物被约束，仿佛被缰绳勒住；另一些则仿佛被刺激，放松缰绳而向前。因为世界，即天主的这辆战车与万物，连天使和星辰自己都引导；它们的运动，虽然多样，却受某些法律约束，我们观看它们在规定给它们自己的时间界限内引导；以致我们现在也情愿与宗徒一同赞叹这位工程师和祂的工程而呼喊：\Quote{啊，天主的智慧和知识的财富多么深奥！祂的审判是多么不可测量，祂的道路是多么不可追寻！}等等。\par

% structure: p0027 source=h2 target=section reason=relative-heading-rank; base=h2

\section{第九章}

\Emph{此外，同样的真理规则教导我们相信，在圣父之后，也相信天主子、耶稣基督我们的主天主，祂就是在旧约中所预许并在新约中所彰显的那一位。}\par

同一真理的规则教导我们，在圣父之后，也要相信天主子、基督耶稣、主我们的天主，但他是那独一无二、即万有创始者之天主的天主子，正如上文已经表明的。因为这位耶稣基督，我愿再次重申，是这一位天主的天主子，我们在旧约中读到祂曾被预许，我们在新约中看到祂显现，以降生成人的真理临在，应验了一切圣事的影子和预象。因为古老的预言和福音都证明祂是亚巴郎之子，达味之子。创世纪本身就预示了祂，当它说：\BibleQuote{我要将它赐给你和你的后裔。}\BibleRef{创 17:8} 当它显示一个人与雅各伯角力时，所说的就是祂；当它说：\BibleQuote{权杖不离犹大，柄杖不离他脚间，直到那应许者来到；祂将是万民的期待。}\BibleRef{创 49:10} 时，也是祂。梅瑟提到祂，当他说：\BibleQuote{请另选派别人去。}\BibleRef{出 4:13} 同一人也再次提到祂，当他作证说：\BibleQuote{天主将从你的兄弟中给你兴起一位先知；你们要如同听从我一样听从祂。}\BibleRef{申 18:15} 他说的也是祂，当他说：\BibleQuote{你的性命将昼夜悬而不定，你必不敢信靠祂。}\BibleRef{申 28:66} 依撒意亚也暗示了祂：\BibleQuote{由叶瑟的根茎将生出枝条，从它的根上将长出花朵。}\BibleRef{依 11:1} 同一人也当他说：\BibleQuote{看，一位贞女将怀孕生子。}\BibleRef{依 7:13} 他论及祂时列举了将从祂而来的治愈，说：\BibleQuote{那时，盲人的眼睛要打开，聋子的耳朵要听见；那时，瘸子要像鹿一样跳跃，哑巴的舌头要说话。}\BibleRef{依 35:3} 他也谈到祂，当他陈明忍耐的德行，说：\BibleQuote{他的声音不闻于街市；压伤的芦苇，他不折断；将熄的灯心，他不吹灭。}\BibleRef{依 13:2} 他也描绘了祂的福音：\BibleQuote{我要与你们订立永久的盟约，即应许于达味的慈爱。}\BibleRef{依 55:3} 他也预言万民将相信祂：\BibleQuote{看，我立祂为万民的领袖和指挥官。不认识的国民要呼求祢，不认识祢的百姓要投奔祢。}\BibleRef{依 55:4} 也是同一人，当论及祂的受难，他呼喊说：\BibleQuote{祂如羊被牵到宰杀之地，又像羔羊在剪毛者前缄默，祂在谦卑中不开口。}\BibleRef{依 53:7} 此外，当他描述祂鞭伤的痕，说：\BibleQuote{因祂的创伤我们得了痊愈。}\BibleRef{依 53:5} 或祂的贬抑：\BibleQuote{我们看见祂，没有俊美，也没有华丽；祂是受苦的人，知道怎样承受病弱。}\BibleRef{依 53:2} 或人民将不信祂：\BibleQuote{我整天向悖逆的百姓伸出双手。}\BibleRef{依 65:2} 或祂将从死者中复活：\BibleQuote{到那日，叶瑟的根子要竖立，祂将兴起统治万民；万民将寄望于祂，祂的安息所将是光荣。}\BibleRef{依 11:10} 或当他论及复活的时间：\BibleQuote{我们将在早晨找到祂，好似为清晨所预备。}\BibleRef{欧 6:3} 或祂将坐在圣父右边：\BibleQuote{上主对我主说：祢坐在我右边，等我使祢的仇敌作祢的脚凳。} 或当祂被描述为万有的拥有者：\BibleQuote{祢向我请求，我必将万民赐祢作产业，将地极赐祢作基业。} 或当祂被显示为万民的审判者：\BibleQuote{天主，求祢将审判权赐给君王，将正义赐给君王的儿子。} 在此我不再继续这个话题：那些关于基督所宣告的事，所有异端者都知道，但那些持守真理的人知道得更清楚。\par

% structure: p0030 source=h2 target=section reason=relative-heading-rank; base=h2

\section{第十章}

\Emph{耶稣基督是天主子，也是真人，与异端者否认祂取了真血肉的幻想相对。}\par

但我要提醒\Emph{你}，基督在福音中的降临，并不应超出那在旧约圣经中由造物主预先所应许的方式；特别是因为那些关于祂的预言已经应验，而那些应验的事早已被预言。正如我有理由可以真实而坚定地对那些异端者的幻象——我不知为何物——说，他们拒绝旧约圣经的权威，把基督说成是从老妇神话中虚构和粉饰出来的：\Quote{你是谁？你从哪里来？谁差遣了你？为什么你现在选择来临？为什么你竟是这样？或者你如何能够来临？或者你为什么不去你所属的地方，除非你表明你并无所有，却来属于他人的？你和造物主的世界有何关系？你和造物主的人有何关系？你和身体的形象有何关系，而你又从中夺去复活的希望？你为什么来觊觎别人的仆人，并想拉拢别人的儿子？你为什么力图把我从主那里夺走？你为什么强迫我亵渎，并对我的父不敬？或者，如果在复活中我失去身体而无法领受自己，那么能从我这里得到什么？如果你愿意拯救，你就该造一个人，以便赐予救恩。如果你渴望救人脱离罪恶，你就该预先赐给我不陷于罪恶的能力。但你随身带着什么法律的认可？你有什么先知话语的见证？或者，我能从你那里期望什么实质的益处，当我看见你来时是幻象而非实体？既然如此，你与身体的形态有何关系，既然你憎恨身体？相反，你在因为憎恨而携带身体实体方面，你将受到驳斥，因为你甚至愿意接受它的形态。因为，如果你憎恨现实，你就应当憎恨对身体的模仿；因为，如果你是不同的存在，你就应当以不同的形态而来，以免你被称为造物主之子，即使你有血肉和身体的相似。诚然，如果你恨恶出生，因为\InnerQuote{恨恶造物主的婚姻结合}，你就应当拒绝任何一个由\InnerQuote{造物主的婚姻}所生的人的相似。}\par

因此，我们既不承认异端者的基督是——如所说的——外观而非实际；因为在他所行的事中，若他本身是幻象而非实际，他就不能行任何实际的事。也不承认那在他身上不带有我们身体任何部分的人，因为\Quote{他从玛利亚那里一无所取}；他也没有来到我们中间，因为他显现为\Quote{异象，而非我们的实体}。我们也不\Emph{承认}那\Emph{是基督}的，即有些人异端所假装的，选择了以太或星质的肉体。我们也不能在他身上看出我们任何救恩，如果我们在他身上甚至不承认我们身体的实体；简言之，也不承认任何其他穿着异端虚构的任何其他种幻躯的人。因为所有这些神话，都被我们主的出生和死亡本身所驳斥。因为若望说：\BibleQuote{圣言成了血肉，住在我们中间}；这样，我们的身体在他内是合理的，因为圣言确实取了我们属于血肉的一切。为此，从他的手脚和肋旁流出血来，为证明他按照我们瓦解的规律死去，分受了我们的身体。而且，他以死时同一个身体实体复活，由那身体本身的伤痕所证明，从而他在自己的肉身中向我们显示了复活的规律，因为他以从我们所得的同一个身体在复活中恢复了。因为复活的规律被确立，即基督以身体的实体复活作为其余人的榜样；因为圣经记载：\BibleQuote{血肉不能继承天主的国}，并非定那由天主亲手建造、不应朽坏的身体实体之罪，而是正当地责斥血肉之罪的过犯，这过犯因人的任意妄为而违犯神圣法律的要求。因为在圣洗圣事和死亡的瓦解中，肉身得以提升并回归救恩，因它被召回无罪的状态，当罪性的死亡被除去时。\par

% structure: p0034 source=h2 target=section reason=relative-heading-rank; base=h2

\section{第十一章}

\Emph{而且，基督不仅是人，也是天主；正如祂是人子，同样祂也是天主之子。}\par

但是，为免我们因宣称我们的主耶稣基督、天主子、造物主是在真身体的\OriginalTerm{实体}{substance}中显现，而似乎同意其他异端者，他们在此主张祂只是且唯独是人，因此想要证明祂是一个赤裸而孤独的人；又为免我们似乎为他们提供了任何反驳的依据，我们并不以这样的方式表达关于祂身体实体的\Emph{教义}，以致说祂只是且唯独是人，而是通过\OriginalTerm{圣言}{Word}的\OriginalTerm{天主性}{divinity}与那物质性的结合，按照圣经坚持祂也是天主。因为说人类救主\Emph{只是人}，这是极大的危险；说万有的主、世界的元首，万物都由父交付给祂、祂父赐予祂一切、祂是万有的所立定、所创造、所安排，万世万代的君王、众天使的元首，在祂面前除父外别无其他者，只是人，并否定祂在这些事中具有天主性的权威。因为异端者的这种轻视也将反作用于天主圣父，如果天主圣父不能生天主子。而且，异端者的任何盲目都不能规定真理。因为他们坚持基督是一样而不坚持另一样，他们看到基督的一面而不看到另一面，也不能因他们所见而夺去我们所未见。因为他们视祂的软弱如同人的软弱，却不算祂的德能如同天主的德能。他们记住肉身的软弱，却排除天主性的德能；然而，如果来自基督软弱的论证有助于证明祂是因软弱而为人，那么从祂的德能中收集的天主性论证也有助于从祂的作为中宣称祂是天主。因为如果祂的苦难显示祂人性的软弱，为什么祂的作为不能宣称祂天主性的德能？因为如果这不能从德能宣称祂是天主，那么苦难也不能从它们显示祂是人。因为无论采取哪一方的原则，都会发现被坚持。所以会有危险：祂不能因祂的苦难而被显示为人，如果祂也不能因祂的德能被证实为天主。那么，我们不可倾向一边而回避另一边，因为任何排除一部分真理的人将永远不能掌握完全的真理。因为圣经既宣布基督也是天主，也同样宣布天主自己是人。圣经既描述耶稣基督是人，也同样描述主基督是天主。因为它不仅说祂是天主子，也说祂是人子；它也不仅说人子，也惯常称祂为天主子。所以，祂既是两者，所以祂是两者，以免祂如果是单独一个，就不能成为另一个。因为正如自然本身规定：凡从人而来者必须相信是为人，同样自然也规定：凡从天主而来者必须相信为天主；但如果祂在从天主而来时不是天主，那么祂即使从人而来，也不能是人。这样，两种教义都会以这种或那种方式受到危害，因为一种被证明失去对另一种的信仰。因此，那些读到耶稣基督人子是人的人，也要读到这同一耶稣也被称为天主和天主子。因为正如作为人，祂出自亚巴郎，同样作为天主，祂在亚巴郎之前就已存在。同样，正如作为人，祂是\Quote{达味之子}，同样作为天主，祂被宣告为达味之主。同样，正如作为人，祂生于\BibleQuote{法律之下}\BibleRef{迦 4:4}，同样作为天主，祂被宣告为安息日的主。同样，正如作为人，祂承受审判，同样作为天主，祂拥有审判活人死人的一切权能。同样，正如作为人，祂生于世界之后，同样作为天主，祂被显示为在世界之前。同样，正如作为人，祂由达味的后裔所生，同样作为天主，世界也被说成是由祂所立定。同样，正如作为人，祂在许多人之后，同样作为天主，祂在万有之前。同样，正如作为人，祂低于他人，同样作为天主，祂比万有更大。同样，正如作为人，祂升到天上，同样作为天主，祂先从天降下。同样，正如作为人，祂往父那里去，同样作为听从父的子，祂将从那里降下。所以，如果祂的软弱证明人性的软弱，祂的威严肯定天主性的德能。因为危险在于，读到两者时，不是相信两者，而是相信其中之一。因此，既然在基督中两者都被读到，就应两者都相信；这样最终信德既是完全的，也是真实的。因为如果在信仰中两个原则中有一个让步，而另一个（而且这正是最不重要的）被拾起作为信仰，真理的准则就被混乱了；这种大胆不会带来救恩，反而因信仰的倾覆而带来死亡的极大危险。\par

% structure: p0037 source=h2 target=section reason=relative-heading-rank; base=h2

\section{第十二章}

\Emph{基督是天主：由旧约圣经的权威证明。}\par

为什么我们犹豫要说圣经不避讳宣告的事呢？为什么信仰的真理要在圣经权威从未犹豫的事上犹豫呢？因为，看，欧瑟亚先知以圣父的口吻说：\Quote{我现今必不借弓、不借马、不借骑兵拯救他们；我却要借主他们的天主拯救他们。}\BibleRef{欧 1:7} 如果天主说祂借天主拯救，而天主除非借着基督并不拯救，那么，人为何要犹豫称基督为天主呢，既然他看到按圣经基督被圣父宣告为天主？是的，如果天主圣父除非借天主并不拯救，那么没人能被天主圣父拯救，除非他宣认基督是天主，圣父许诺在他内并藉他赐予救恩；这样，有理由地，凡承认他是天主的，可在基督天主内找到救恩；凡不承认他是天主的，将失去不能在他处找到的救恩，除非在基督天主内。因为正如依撒意亚说：\Quote{看，一位贞女将怀孕生子，你要给他起名叫\OriginalTerm{厄玛奴耳}{Emmanuel}，即译为：\OriginalTerm{天主与我们同在}{God with us}。}\BibleRef{依 7:14} 同样，基督自己说：\Quote{看，我同你们天天在一起，直到今世的终结。}\BibleRef{玛 28:20} 所以，他是天主与我们同在；是的，更是，他在我们内。基督与我们同在，因此，他的名字是天主与我们同在，因为他也与我们同在；或者难道他不与我们同在？那么他如何说他在我们当中？那么，他确实与我们同在。但因为他与我们同在，他被称为厄玛奴耳，即天主与我们同在。天主，因此，因为他与我们同在，被称为天主与我们同在。同一位先知说：\Quote{无力的手，你们要坚强！无力的膝，你们要稳健！心怯的人，你们要鼓起勇气，要坚强，不要害怕！看，我们的天主要来施行审判；他要亲自来拯救你们。那时，盲人的眼要睁开，聋子的耳要开启；那时，瘸子要跳跃如鹿，哑巴的舌头要欢呼。} 先知说在天主来临时，这些记号将要发生；让人们承认，要么基督是天主子，在他的来临并藉着他这些医治的奇迹成就了；要么，被基督天主性的真理所折服，让他们冲入另一个异端，拒绝承认基督是天主子、是天主，而宣称他是圣父。因为，被先知的话所束缚，他们不能再否认基督是天主。那么，当那些记号称要在天主来临之际发生，而这些记号在基督来临时已显明，他们如何回答？他们如何接受基督为天主？因为他们现在不能否认他是天主。作为天主圣父还是作为天主圣子？如果作为圣子，为何他们否认天主子是天主？如果作为圣父，为何他们不追随那些似乎坚持此类亵渎的人？除非因为在我们与他们的这场关于真理的辩论中，目前对我们来说这就够了，即他们以任何方式被说服，承认基督是天主，因为他们甚至曾想否认他是天主。他借哈巴谷先知说：\Quote{天主将从南方而来，圣者从黑暗而稠密的山而来。} 他们希望谁是从南方来的？如果说那是全能的天主圣父，那么天主圣父来自一个地方，从那个地方，他就这样被排除，并被限制在某住所的狭窄之内；这样，正如我们所说，\OriginalTerm{撒伯流}{Sabellius}的亵渎异端就被这些人体现出来。因为基督不被相信是圣子，而是圣父；既然他们声称他严格地只是一个\OriginalTerm{纯粹的人}{bare man}，以一种新的方式，基督又被这些人证明是全能的天主圣父。但如果基督生于白冷，该地区的地理分区朝向天空的南方部分，他由圣经也被称为天主，这位天主被合适地描述为从南方而来，因为他被预见将从白冷而来。那么，让他们从两个选择中选出他们更喜欢的一个，那个从南方来的是圣子，还是圣父；因为圣经说天主将从南方来。如果是圣子，为何他们不敢称他为基督和天主？因为圣经说天主要来。如果是圣父，为何他们不敢与撒伯流的冒失相连，后者说基督是圣父？除非因为，无论他们称他为圣父或圣子，从他的异端，尽管不情愿，他们必须退出，如果他们习惯于说基督只是人；当被事实本身所迫，他们将快要高举他为天主，无论他们希望称他为圣父或希望称他为圣子。\par

% structure: p0040 source=h2 target=section reason=relative-heading-rank; base=h2

\section{第十三章}

\Emph{这同一真理由新约的圣经得以证明。}\par

同样，若望描述基督的诞生，说：\Quote{圣言成了血肉，居住在我们中间，我们见了祂的光荣，正如父独生者的光荣，满溢恩宠和真理。}。再者，\Quote{祂的名字称为天主的圣言，} \BibleRef{默 19:13}，并非毫无理由。我心发出美好的言语；祂随后推论性地用君王的名字称呼这言语，\Quote{我要向君王陈述我的事工。} 因为\Quote{藉着祂万有都是造成的，没有祂什么也没有造成。} \BibleRef{若 1:3} \Quote{无论是}宗徒说\Quote{座天使或宰制者或掌权者或异能者，有形无形的万有，都藉着祂而存在。} \BibleRef{哥 1:16} 此外，这就是来到自己地方的圣言，自己的人却没有接受祂。因为世界是藉着祂造的，世界却不认识祂。\BibleRef{若 1:10} 此外，这圣言\Quote{在起初就与天主同在，圣言就是天主。} \BibleRef{若 1:1} 那么，当最后一句说\Quote{圣言成了血肉，居住在我们中间}时，谁能怀疑基督——祂的诞生，且因祂成了血肉，而是人；又因祂是天主的圣言——谁敢迟疑而不毫不犹豫地宣告祂是天主呢？尤其是当他思量福音书，它已将这双重实体结合在基督诞生的一同和谐之中。因为正是祂，\Quote{像新郎走出洞房，又如壮士欢跃奔路。祂的出发从天涯，祂的回归到地极。} 因为，即使到至高之处，\Quote{除了从天降下的人子，没有谁升过天。} \BibleRef{若 3:13} 重复同样的话，祂说：\Quote{父啊，求祢以我在创世之前与祢同在时所拥有的光荣，光荣我吧。} 如果这圣言如新郎降临到血肉，以便通过接纳血肉，祂能作为人子升到那里，而天主圣子曾作为圣言从那里降下；那么，由于相互的结合，血肉佩戴着天主的圣言，天主圣子承载着血肉的脆弱；当血肉被许配，升到那里——它曾无血肉地从那里降下——最终获得那光荣，这光荣被显示为在创世之前就已拥有，从而最明显地证明祂是天主。然而，虽然这世界本身被说成是在祂之后建立的，却发现是祂创造的；藉着祂内的那同一神性——世界由此被造成——祂的光荣和权柄都得到证明。此外，如果说只有天主才能洞察人心的隐秘，而基督洞察人心的隐秘；如果说只有天主才能赦罪，而同一基督赦罪；如果说没有人能从天上降下，而祂却从天降下；如果说这句话对任何人都不能成立：\Quote{我与父原是一体，} \BibleRef{若 10:30} 而唯独基督出于对祂神性的意识宣告了这句话；最后，如果宗徒多默，受教于关于基督神性的一切证据和条件，回答基督说\Quote{我主！我天主！} \BibleRef{若 20:28}；此外，宗徒保禄说\Quote{圣祖也是他们的，并且基督按血统说也是从他们来的，祂是万有之上，永远受赞美的天主。} \BibleRef{罗 9:5}，他在书信中写道；并且同一宗徒宣称他被立为\Quote{宗徒，不是由于人，也不是藉着人，而是藉着耶稣基督；} 并声称他学得福音不是从人或藉着人，而是从耶稣基督领受的；那么，基督是天主，这是合理的。因此，在这方面，必须确立两件事之一。因为既然万有都是藉着基督造成的，祂要么在万有之先，因为万有藉着祂；这样祂就是天主；要么因为祂是人，祂就在万有之后，那么万有就都不是藉着祂造成的。但我们不能说万有不是藉着祂造成的，因为我们看到圣经记载万有都是藉着祂造成的。因此祂不在万有之后；也就是说，祂不仅是人——人是在万有之后——也是天主，因为天主在万有之先。因为祂在万有之先，因为万有都是藉着祂；而如果祂只是人，那么万有就不是藉着祂；或者如果万有是藉着祂，祂就不只是人，因为如果祂只是人，万有就不是藉着祂；不，万有就都不是藉着祂。那么，他们如何回答呢？说万有不是藉着祂，所以祂只是人？那么万有如何是藉着祂呢？因此祂不只是人，也是天主，因为万有是藉着祂；所以我们应当合理地理解，基督不仅是人——人在万有之后——也是天主，因为万有都是藉着祂造成的。因为你怎么能说祂只是人，而你看见祂也在肉身中呢？除非当两方面都被考虑时，两个真理都被正确相信。\par

% structure: p0043 source=h2 target=section reason=relative-heading-rank; base=h2

\section{第十四章}

\Emph{作者继续同一论证。}\par

然而，异端者仍然不愿承认基督是天主，尽管他已从如此多的言语和事实中看到基督被证明是天主。如果基督只是人，那么当他来到这个世界时，他怎能来到属于自己的地方呢？因为人不能创造世界。如果基督只是人，那么世界怎能说是藉着他造的？因为世界不是由人造成的，而是人是在世界之后被创造的。如果基督只是人，那么基督怎会不但是达味的后裔，而且还是圣言成了肉身并住在我们中间呢？因为虽然原祖不是从种子生的，但原祖也不是由圣言和肉身的结合而成的。因为那圣言并不是成了肉身，也没有住在我们内。如果基督只是人，那么\BibleQuote{从天上来的，见证他所见所闻的}\BibleRef{若 3:31}又怎会成立呢？因为很明显，人不能从天上而来，因为他不能在那里出生。如果基督只是人，那么\BibleQuote{看得见的事物和看不见的事物，上座者、宰制者、率领者}怎能说是藉着他并在他内受造的呢？因为天上的能力不能由人造成，因为它们在人之先。如果基督只是人，那么他怎会在他被呼求的任何地方临在呢？因为能临在于每一地方的不是人的本性，而是天主的本性。如果基督只是人，那么为什么在祈祷中呼求一个人作为中保呢？因为呼求一个人来提供救恩被认为无效。如果基督只是人，那么为什么寄望于他呢？因为寄望于人被宣告为可咒骂的。如果基督只是人，那么为什么否认基督不会导致灵魂的毁灭呢？因为据说得罪人可以被宽恕。如果基督只是人，那么若翰洗者怎会见证并说：\BibleQuote{那在我以后来的，成了在我以前的，因为他在我以前}\BibleRef{若 1:15}？如果基督只是人，那么\BibleQuote{圣父所做的事，圣子也照样做}\BibleRef{若 5:19}怎么成立呢？因为人不能做与天主天上的操作相似的工作。如果基督只是人，那么\BibleQuote{就如圣父在自己内有生命，同样也赐给圣子在自己内有生命}\BibleRef{若 5:26}怎么成立呢？因为人不能像圣父那样在自己内有生命，因为他不是在永世中荣耀的，而是由可朽的材料造的。如果基督只是人，那么他怎能说：\BibleQuote{我是从天上降下来的永生的食粮}\BibleRef{若 6:51}？因为人既不能成为生命之粮（他自己是可朽的），也不能从天上降下来，因为没有任何可朽之物安置在天上。如果基督只是人，那么他怎能说：\BibleQuote{从来没有人见过天主，除了那出于天主的，他见过天主？}因为如果基督只是人，他不能看见天主，因为没有人见过天主；但如果他是出于天主的，他见过天主，他就希望人理解他超过人，因为他见过天主。如果基督只是人，那么他为什么说：\BibleQuote{如果你们看见人子升到他以前所在之处，那会怎样呢？}但他升了天，所以他曾在那里，因为他回到他以前所在之处。但如果他被圣父从天上派遣而来，他肯定不只是人；因为如我们所说，人不能从天上而来。因此，作为人他以前不在那里，而是升到了他不在的地方。但天主的圣言——我所说的天主的圣言，以及天主，万物是藉着他造的，没有他什么也没有造——降下来了。因此，不是人从天上而来，而是天主的圣言；也就是说，天主从那里降下来了。\par

% structure: p0046 source=h2 target=section reason=relative-heading-rank; base=h2

\section{第十五章}

\Emph{他又从福音证明基督是天主。}\par

如果基督只是人，那么祂说：\BibleQuote{我虽然为自己作证，但我的见证是真实的，因为我知道我从哪里来，往哪里去；你们却不知道我从哪里来，往哪里去。你们是按肉身判断的？}这是怎么回事？看，祂还说，祂要回到祂以前作证说祂从那里来的地方去——也就是从天上——因为祂是被派遣的。因此，祂从祂所来的地方降下，同样也要回到祂所降下的地方去。所以，如果基督只是人，祂就不会从那里来，因此也不会去那里，因为祂不会是从那里来的。此外，借着从那里来——而作为人祂不可能从那里来——祂表明了自己是以天主的身份来的。因为那些愚昧无知、对这降临之事毫无了解的犹太人，使这些异端者成了他们的后继者，因为对他们说了：\BibleQuote{你们不知道我从哪里来，往哪里去；你们是按肉身判断的。}他们和犹太人一样，只认为基督有肉身的出生，相信基督不过是一个人；没有考虑这一点：既然人不能从天上下来，也不能回到那里去，那么从那里降下来的必定是天主，因为人不可能从那里来。如果基督只是人，祂怎么说：\BibleQuote{你们是出于下，我是出于上；你们是这世界的，我不是这世界的？}但因此，如果每一个人都属于这个世界，而基督因此也在这世界里，祂只是人吗？决不可以！但请思考祂说的话：\BibleQuote{我不是这世界的。}那么如果祂只是人，祂说\BibleQuote{这世界的}岂不是说谎？或者，如果祂不说谎，祂就不是这世界的；所以祂不仅是人，因为祂不是这世界的。但为了不让祂是谁成为秘密，祂宣告了祂是从哪里来的：\BibleQuote{我，}祂说，\BibleQuote{是从上来的}，也就是说，从天上来的，人不能从天上来，因为人不是在天上造的。因此，祂是从上来的天主，所以祂不属于这世界；虽然，在某种程度上，祂也属于这世界：因此基督不仅是天主，也是人。正如按照祂作为圣言的天主性而言，祂不属于这世界；同样，按照祂所取的肉身的脆弱性而言，祂属于这世界。因为人与天主结合，天主与人相连。但基督在这里更强调祂独一的天主性这一方面，因为犹太人的盲目只看基督的肉身这一方面；因此，在当前的段落中，祂略过属于这世界的身体的脆弱性，只谈祂不属于这世界的天主性：这样，既然他们倾向于相信祂只是人，基督就越发引导他们思考祂的天主性，以便相信祂是天主，祂渴望克服他们对祂天主性的不信，暂时不提祂的人性状况，只将祂的天主性摆在他们面前。如果基督只是人，祂怎么说：\BibleQuote{我出于天主而来到}，而显然人是天主所造的，并非出于天主？但正如作为人祂不出于天主，同样，天主的圣言却是出于天主的，关于这圣言经上说：\BibleQuote{我的心中涌出美好的言语}；这圣言既然出于天主，也合理地与天主同在。而且，因为它并非没有效果而发出，所以它合理地造了万有：\BibleQuote{万物是借着祂造成的，凡受造的，没有一样不是借着祂造成的。}而造万物的圣言（就是天主）。\BibleQuote{天主，}经上说，\BibleQuote{就是圣言。}因此，天主出于天主，因为出于的圣言是天主，祂出于天主。如果基督只是人，祂怎么说：\BibleQuote{谁若遵守我的话，就永远见不到死亡？}永远见不到死亡！这不就是不死吗？但不死是与天主性相关联的，因为天主性是不死的，而不死是天主性的果实。因为每一个人都是可死的；不死不可能从那可死的而来。因此，从作为可死之人的基督那里，不可能产生不死。\BibleQuote{但是，}祂说，\BibleQuote{谁遵守我的话，就永远见不到死亡}；因此基督的话赋予不死，借着不死赋予天主性。虽然不能认为一个自己可死的人能使另一个人不死，但基督的话不仅宣告，而且实际赋予不死：显然，祂不只是人，因为赐予不死，如果祂只是人就不能赐予；但借着不死赐予天主性，祂证明自己就是天主，因为赐予天主性，如果祂不是天主就不能赐予。如果基督只是人，祂怎么说：\BibleQuote{在亚伯拉罕以前，我就有？}因为没有人能在他自己由之而来者之前存在；也不可能有人先于他自身起源的那一位。然而基督虽然生于亚伯拉罕，却说祂在亚伯拉罕以前。因此，要么祂说的不是真的，并且在欺骗——如果祂不在亚伯拉罕以前，因为祂出于亚伯拉罕；要么祂不欺骗——如果祂也是天主，并且在亚伯拉罕以前。若非如此，既然出于亚伯拉罕，祂就不可能在亚伯拉罕以前。如果基督只是人，祂怎么说：\BibleQuote{我认识他们，我的羊跟随我；我赐给他们永生，他们永不丧亡？}然而，既然每一个人都受制于可死的规律，因此不能永久保守自己，就更不能永久保守另一个人了。但基督许诺永远赐予救恩，如果祂不赐予，祂就是欺骗者；如果祂赐予，祂就是天主。但祂不欺骗，因为祂赐予祂所许诺的。因此，祂是天主，祂赐予永久的救恩，而人自己不能\Emph{永久}保守自己，也不能赐予他人。如果基督只是人，祂说：\BibleQuote{我与圣父原是一体}是什么意思？因为\BibleQuote{我与圣父原是一体}如何可能，如果祂不是天主又是圣子呢？——因此祂可以被称为一体，因为祂是出于祂自身，既是祂的圣子，又由祂而生，又被宣告出于祂，由此祂也是天主；当犹太人认为这可憎，并以为是亵渎，因为祂在这些言论中表明自己是天主，于是立刻冲向祂，要拿石头打祂，并狂热地投掷石头，祂借着圣经的实例和见证有力地反驳了反对者。\BibleQuote{如果，}祂说，\BibleQuote{那些承受天主话的人，祂尚且称他们为神，而圣经是不能废去的，那么圣父所祝圣并派遣到世界上来的，你们却说：你说亵渎的话，因为我曾说：我是天主子。}这些话中，祂并没有否认自己是天主，反而更确认了祂是天主的断言。因为毫无疑问，那些承受天主话的人被称为神，而那位超越所有这些的，更是天主。然而，祂以合宜的方式，用合法的技巧恰当地驳斥了亵渎的诽谤。因为祂希望这样被理解：祂是天主，但作为天主子，而不愿被理解为圣父本身。因此祂说祂是\BibleQuote{被派遣的}，并表明祂从圣父那里显了许多善工；由此祂希望不要被理解为圣父，而是圣子。在祂辩护的后半部分，祂提到了圣子，而不是圣父，当祂说：\BibleQuote{你们说：你说亵渎的话，因为我曾说：我是天主子。}因此，就亵渎的罪责而论，祂称自己为圣子，不是圣父；但就祂的天主性而论，祂说：\BibleQuote{我与圣父原是一体}，证明了自己是天主子。所以祂是天主，但作为是天主，祂是圣子，而不是圣父。\par

% structure: p0049 source=h2 target=section reason=relative-heading-rank; base=h2

\section{第十六章}

\Emph{再以福音证明基督是天主。}\par

如果基督仅是凡人，祂自己怎么说：\Quote{凡信我的，必永远不死？}然而，独自信靠人的人被称为可咒骂的；但信基督的人并不被咒骂，反而被说成必永远不死。因此，如果一方面祂仅是凡人，如异端者所主张的那样，那么任何信祂的人怎么会永远不死呢？因为信赖人的人被视为可咒骂的。另一方面，如果他不被咒骂，反而如经上所读，注定获得永生，那么基督不仅是凡人，也是天主；信祂的人既免除了咒诅的一切危险，也达到了义德的果实。如果基督仅是凡人，祂怎么说圣神\Quote{要从我那里领受，并将那些事宣示给你们？}因为圣神不从人领受任何事物，而是圣神将知识赐予人；圣神也不从人学习未来之事，而是教导人关于未来。因此，要么圣神没有从作为人的基督领受祂该宣示的，因为人不能给圣神任何东西，既然人自己应从圣神领受，而基督在此事上既错误又欺骗，说圣神要从身为人的祂领受祂所宣示的；要么祂不欺骗我们——事实上祂并没有——而圣神已从基督领受了祂所宣示的。但若圣神已从基督领受了祂要向我们宣示的，则基督大于圣神，因为圣神若不是小于基督，就不会从基督领受。然而圣神小于基督，此外，这一点本身证明基督是天主，因为圣神从祂领受了所宣示的：因此，基督天主性的见证是巨大的，因为在这项安排中，圣神被发现小于基督，并从祂领受而赐予别人；因为如果基督仅是凡人，基督就会从圣神领受祂该说的话，而不是圣神从基督领受祂该宣示的。如果基督仅是凡人，祂为什么为我们立下这样的信仰规则，在其中祂说：\BibleQuote{永生就是：认识你，唯一的真天主，和你所派遣来的耶稣基督。}如果祂不希望自己也被理解为天主，为什么祂加上\Quote{和你所派遣来的耶稣基督}，除非是因为祂也希望被接受为天主？因为如果祂不希望被理解为天主，祂就会加上\Quote{和你所派遣来的那人耶稣基督}；但事实上，祂并没有加上这个，基督也没有将自己交付给我们仅作为人，而是将自己与天主联合，正如祂希望通过这种结合使自己也被理解为天主，如同祂是的那样。因此，我们必须按照所规定的规则，相信主、唯一真天主，因此也相信祂所派遣的耶稣基督；祂绝不会，如我们所说，将自己与圣父联结，除非祂希望自己也被理解为天主：因为如果祂不希望被理解为天主，祂就会将自己与祂分开。如果祂知道自己仅是凡人，祂就会将自己仅置于人当中；如果祂不知道自己是天主，祂就不会将自己与天主联结。但在这种情况中，祂对自己的为人保持沉默，因为没有人怀疑祂是人，并有理由将自己与天主联结，以便为那些应相信的人确立祂天主性的规范。如果基督仅是凡人，祂怎么说：\BibleQuote{现今，以我未有世界以前同你所享有的光荣，光荣我。}？如果在未有世界以前，祂与天主同享光荣，并保持祂与圣父同享的光荣，那么祂在未有世界以前就已存在，因为除非祂自己先存在，能够保持那光荣，祂就不会拥有那光荣。因为无人能拥有任何事物，除非他自己先存在才能保持任何事物。但现在基督在创立世界以前就有光荣；因此祂自己是在创立世界以前就存在的。因为除非祂在创立世界以前就已存在，祂就不能在创立世界以前拥有光荣，因为祂自己当时并不存在。但人的确不能在创立世界以前有光荣，因为人是在世界之后；但基督却有——因此祂是在世界之先。因此祂不仅是人，因为祂在世界之先。因此祂是天主，因为祂在世界之先，并在世界之前拥有祂的光荣。也不要让这以预定来解释，因为这样表达并不如此。或者让那些这样想的人加上这一点，但增删圣经的人有祸了，正如删减圣经的人一样。因此，那不可说的，就不可加添。因此，预定被排除，既然它没有这样定规，基督在本体上是在创立世界以前。因为祂是\BibleQuote{圣言，万物藉祂而造，没有祂什么也没有造。}因为即使祂被说成在预定中光荣，且这预定是在创立世界以前，秩序仍须保持，并且有许多人是在祂之前被预定得光荣的。因为就那预定而言，如果基督在被他们之后被指定，那么祂将被视为小于其他的人。因为如果这光荣是在预定中，基督最后才接受了光荣的预定；因为在祂之前，\Person{亚当}、\Person{亚伯尔}、\Person{哈诺客}、\Person{诺厄}、\Person{亚巴郎}和许多其他人都已被预定。因为既然在天主那里，所有的人和事物的秩序都已安排，许多人在基督获光荣的预定之前就被说成是预定的。这样，基督被发现是次于其他人，即使事实上祂比天使本身更好、更大、更古老。那么，要么让所有这些被搁置一边，以便基督的天主性被摧毁；要么如果这些不能被搁置，就请异端者将祂应有的天主性归于基督。\par

% structure: p0052 source=h2 target=section reason=relative-heading-rank; base=h2

\section{第十七章}

\Emph{此外，\Person{梅瑟}在\WorkTitle{圣经}开端也证明了这一点。}\par

倘若\Person{梅瑟}遵循这同一真理的法则，并在其神圣著作的开端向我们传授这一原则，使我们得以知晓万有都是由天主子、即\OriginalTerm{圣言}{Word of God}所创造和完成的？因为他说了与\Person{若望}和其他人所说的相同的话；不仅如此，\Person{若望}和其他人明显是从他那里领受了他们所说的话。若\Person{若望}说：\BibleQuote{万有是靠祂而造成的；凡受造的，没有一样不是由祂而造成的}，先知\Person{达味}也说：\BibleQuote{我向君王陈述我的事工}。此外，\Person{梅瑟}描述天主起初命令要有光，要建立穹苍，要使水聚在一处，要使旱地出现，要使地生出果实各按其种，要生出动物，要在天上设立光体和星辰。他表明当时没有其他一位临在于天主——这些工程由祂命令而作成——除了那一位万有都是藉祂而造成的、没有祂什么也不能造成的主。如果祂是天主的圣言——\BibleQuote{因为我的心涌出美言}——他便显明在起初已有圣言，圣言与圣父同在，而且圣言就是天主，万有是藉祂而造成的。此外，这\BibleQuote{圣言成了血肉，居住在我们中间}——即天主子基督；我们后来按照肉身接受祂为人，并看到祂在创世之前就是天主的圣言和天主，我们依照旧约和新约的教导，合理地相信和持守祂既是天主又是人，即基督耶稣。倘若同一位\Person{梅瑟}描述天主说：\BibleQuote{让我们按我们的肖像和模样造人}；其后又说：\BibleQuote{天主于是造了人；按天主的肖像造了他，造了他们男和女}？如果，如我们已经证明的，万有是由天主子所造成的，那么人类也当然是由天主子所制定的，万有是为了人类而造成的。再者，当天主命令造人时，那位造人的被称为天主；但天主子造了人，也就是说，天主的圣言，\BibleQuote{万有是靠祂而造成的，凡受造的，没有一样不是由祂而造成的}。这圣言成了血肉，居住在我们中间；因此基督是天主；因此人是由基督作为天主子造成的。但天主按天主的肖像造了人；因此那位按天主的肖像造人的就是天主；因此基督是天主：这样，旧约关于基督位格的见证有理由不摇摆，因为它得到了新约彰显的支持；新约的权能也不受减损，因为它的真理扎根于同一旧约的根源。因此，那些妄断天主子基督既是天主又是人，却以为祂只是人而非天主的人，是违背旧约和新约，因为他们败坏了旧约和新约的权威与真理。倘若同一位\Person{梅瑟}处处描述天主圣父是无限无量、没有终结的，并非被任何地方所围困，而是那包罗万有的；祂不在一个地方，而是万有都在祂内，祂统御万有、拥抱万有，因此祂既不能下降也不能上升，因为祂自己统御并充满万有；然而，他却描述天主下降要察看世人所建造的塔，并问说：\BibleQuote{来吧；让我们下去，在那里混乱他们的语言，叫他们彼此听不懂对方的话}。他们在此声称那位下降到那座塔、并要来察看那些人的天主是谁？是天主圣父吗？那么祂就受地方限制；祂又如何拥抱万有？或者说是一位天使与天使一起下降，说：\BibleQuote{来吧}；随后\BibleQuote{让我们下去，在那里混乱他们的语言}？然而在《申命纪》中我们看到天主说了这些事，写着说：\BibleQuote{至高者将万民分开时，按天主的众天使的数目，划定了各民族的疆界}。因此，圣父既没有下降，正如主题本身所示；天使也没有命令这些事，正如事实所显明。那么，必然是从天降下的那位，宗徒\Person{保禄}说：\BibleQuote{那下降的，正是那升到诸天之上的，为要充满万有}，即天主子，天主的圣言。但天主的圣言成了血肉，居住在我们中间。这就是基督。因此必须宣认基督是天主。\par

% structure: p0055 source=h2 target=section reason=relative-heading-rank; base=h2

\section{第十八章}

\Emph{此外，从\Person{亚巴郎}所见的那位被称为天主这一事实；这不能理解为圣父，因为从来没有人见过祂，而是理解为取了天使形象的\OriginalTerm{圣子}{Son}。}\par

看，同一位\Person{梅瑟}在另一处告诉我们\BibleQuote{天主被亚巴郎看见了}。然而同一位\Person{梅瑟}从天主那里听到\BibleQuote{没有人看见天主而能存活}。如果天主不能看见，如何被看见？或者如果祂被看见了，怎么又不能被看见？因为\Person{若望}也说：\BibleQuote{从来没有人见过天主}；宗徒\Person{保禄}说：\BibleQuote{从来没有人看见过祂，也不能看见祂}。但圣经当然不会说谎；因此，的的确确，天主被看见了。由此可以明白，被看见的不是圣父，因为祂从未被看见；而是圣子，祂习惯于下降，并且因下降而被看见。因为祂是不可见的天主的肖像，正如人类状况的不完美和脆弱，有时甚至那时也习惯于在\OriginalTerm{不可见的天主的肖像}{image of the invisible God}中看见天主圣父，即在\OriginalTerm{天主子}{Son of God}中。因为人类的软弱要逐步、循序渐进地通过肖像得到强化，以达到将来有一天能看见天主圣父的荣耀。因为伟大的事物若突然而来，是危险的。因为黑暗之后的突然阳光，以其过度的光辉，不会使不适应的眼睛看清白天的亮光，反而会使它们失明。\par

为避免此事损伤人眼，黑暗便逐渐被驱散；那光体以微小而不知不觉的增升升起，借其光芒的\Emph{柔和}增强，温柔地使人的眼睛习惯于承受它的全圆。因此，基督——即天主的肖像、天主之子——也照人所能见的样子被人观看。如此，人类命运中的软弱与不完美由祂滋养、引领并教育，使人习惯于瞻仰圣子，终有一日能得见天主圣父本身，如祂之所是，免得被祂骤然难忍的辉煌所击，以至于无法看见始终渴望的天主圣父。因此，所见的是圣子；而天主之子是天主圣言；圣言成了血肉，居住在我们中间；这就是基督。世间有什么理由让我们犹豫不称祂为天主？祂以如此多的方式被承认被证明是天主。此外，天使遇见撒辣的使女哈加尔，她在家乡被逐且被赶走，在往苏尔路上的水泉旁；天使询问并得知她逃跑的原因，随后劝她自谦；又赐她为母的盼望，应许并保证从她胎中要生出众多后裔，且她将生下依市玛耳；又述说他居住的地方，描述他的生活方式；然而圣经称此天使为主和天主——因为若非这天使也是天主，祂不会应许后裔的祝福。让异端者看看他们如何理解此段。那被哈加尔看见的是圣父吗？因为祂被称为天主。但我们绝不可称天主圣父为天使，免得祂从属于另一位，成为其天使。但他们会说那是天使。那么祂如何又是天主呢？因为这一称号从未赐予天使，除非双方真理迫使我们持此见解：我们应当明白那是天主圣子，祂因出自天主而被称为天主，因为祂是天主之子。但因祂服从圣父、宣告圣父的旨意，祂被称为\OriginalTerm{大谋议天使}{Angel of Great Counsel}。因此，虽然此段既不适合圣父（免得祂被称为天使），也不适合天使（免得他被称为天主），却适合基督：祂是天主，因为祂是天主之子；祂也是天使，因为祂宣告圣父的意念。异端者应当明白他们是在反对圣经，因为当他们说相信基督也曾是天使时，却不愿宣告祂也是天主，尽管他们在旧约中读到祂常来探访人类。此外，摩西又加上天主在玛默勒橡树旁显现给亚巴郎，那时他正坐在帐幕门口，在正午。然而，虽然他看见了三个人，\Emph{请注意}他称其中一位为主；他洗了他们的脚，献上灰烬烤的饼、奶油和大量牛奶，并央请他们留下作客进食。之后他听见自己将成为父亲，得知他的妻子撒辣将为他生一个儿子；又得知索多玛人民因当受的惩罚而被毁灭，并明白天主因索多玛的呼声而降临。在此处，如果他们坚持认为那时圣父曾与两位天使一同被接待，那么异端者便相信圣父是可看见的。但如果是一位天使，尽管三位天使中有一位被称为主，为何——虽然不常见——天使被称为天主呢？除非，为恢复圣父的不可见性，并为天使保留其从属地位，只有天主圣子——祂也是天主——被亚巴郎看见，并被相信曾受接待。因为祂以圣事方式预表了祂将来要成为的。祂作了亚巴郎的客人，将要成为亚巴郎子孙中的一员。祂洗了他们的脚，以此证明祂是谁，并在祂身上归还在从前由圣父所行的接待恩情。因此，为叫人无疑祂就是毁灭索多玛人民时作亚巴郎客人的那一位，经上宣告：\BibleQuote{那时，主从天上从主那里降下硫磺与火，落在索多玛和哈摩辣身上。}因为先知也以天主的口气说：\BibleQuote{我颠覆了你们，如同主颠覆了索多玛和哈摩辣。}因此，主颠覆了索多玛，即天主颠覆了索多玛；但在颠覆索多玛时，主从主那里降下火。这主就是亚巴郎所见的天主；这天主就是亚巴郎的客人，当然可见，因为祂也是可触摸的。虽然圣父是不可见的，当时确实未被看见，但那位习惯于被触摸和被看见的，被看见并受接待。这就是天主圣子，\BibleQuote{主从主那里降下硫磺与火，落在索多玛和哈摩辣身上。}这就是天主圣言。圣言成了血肉，居住在我们中间；这就是基督。因此，作亚巴郎客人的不是圣父，而是基督。那时被看见的也不是圣父，而是圣子；基督被看见了。所以，基督正确被称为主和天主，祂被亚巴郎看见，正是因为祂作为天主圣言，在亚巴郎之前已由天主圣父所生。此外，圣经说，同一位天使和天主探访并安慰同一位哈加尔，当时她同儿子被赶出亚巴郎的住所。因为当她在旷野把婴儿放下，皮袋里的水用尽了，孩子喊叫，她举哀哭泣，\BibleQuote{天主听见了孩子从所在之处发出的声音。}在说明是天主听见了婴儿的声音后，又加上：\BibleQuote{天主的使者从天上呼唤哈加尔，}称那位曾经被称为天主的为使者，并宣告那位曾显为使者的为主；这位天使和天主又向哈加尔本人应许更大的安慰，说：\BibleQuote{不要怕，因为我已听见孩子从所在之处发出的声音。起来，抱起孩子，用手扶持他，因为我要使他成为大国。}为什么这位使者，如果只是天使，能自称拥有说\BibleQuote{我要使他成为大国}的权柄？因为这类权柄显然属于天主，不能属于天使。因此，祂被证实为天主，因为祂能行这事；为证明这一点，圣经立即补充说：\BibleQuote{天主开了她的眼，她看见一口活水井；她去了，从井里灌满皮袋，给那孩子喝；天主与那孩子同在。}那么，如果这位天主与主同在，开了哈加尔的眼使她看见活水井，并因孩子口渴的急需汲水，而这位从天上呼唤她的天主被称为天使——但之前听见孩子哭喊的天主其实是天主——那么，祂不能被理解为不是天使，正如祂也是天主一样。既然这不适用于或不适合圣父——祂只是天主——却适用于基督，祂被宣告不仅是天主，也是天使，那么显然，对哈加尔说话的不是圣父，而是基督，因为祂是天主；天使的名号也归于祂，因为祂成了\OriginalTerm{大谋议天使}{Angel of Great Counsel}。祂是天使，因为祂揭示了圣父的怀抱，正如若望所陈述的。因为若望自己说，那位揭示圣父的怀抱的，作为圣言，成了血肉以揭示圣父的怀抱，那么基督不仅是人，也是天使；不仅天使，圣经也显示祂是天主。我们相信这是真实的；因此，如果我们不同意理解为那时对哈加尔说话的是基督，我们就必须要么使天使成为天主，要么将全能的天主圣父算在天使之中。\par

% structure: p0059 source=h2 target=section reason=relative-heading-rank; base=h2

\section{第十九章。}

\Emph{天主也曾以天使的形象显现于雅各伯；即天主子。}\par

如果我们在别处也读到同样描述天主为天使，该如何？因为当雅各伯向他的妻子肋阿和辣黑耳抱怨她们父亲的不义，并告诉她们他现在想要离开返回本地时，他还引用了梦的权威；这时他说天主的使者在梦中曾向他说：\Quote{雅各伯、雅各伯。我说，}他说，\Quote{是什么？举目观看，公山羊和公绵羊跳在羊群上，母山羊是黑色、白色、杂色、灰白和有斑点的：因为我已看见拉班对你所做的一切。我是天主，曾在天主的地方向你显现，你在那里为我膏抹了石柱，并向我许愿。现在起来，离开这地，往你本乡去，我必与你同在。}如果天主的使者这样对雅各伯说话，且使者自己提及并说：\Quote{我是天主，曾在天主的地方向你显现}，我们毫不迟疑地看到，这不仅被宣告为天使，也是天主；因为祂说到雅各伯在天主的地方向祂本人所许的愿，祂没有说\Emph{在我的地方}。因此那是天主的地方，祂也是天主。此外，经上只写\Quote{在天主的地方}，没有说\Quote{在天使和天主的地方}，而只是\Quote{天主的地方}；那应许这些事的那位被显明既是天主也是天使，因此合理地，那位只被称为天主的，与那位不仅被称为天主、也被称为天使的，必须有区别。因此，若如此大的权威不能归于任何别的天使，以至于祂应自称为天主，并见证有愿向祂许过，除非唯独归于基督——向祂不仅可以作为天使，也可以作为天主许愿；那么显然，祂不应被理解为父，而应被理解为子、天主和天使。此外，如果这是基督——确实如此——那么说基督只是人或只是天使、扣留神圣名号的权能的人，就处于可怕的危险中；这权能是祂常凭天上圣经的信德所领受的，这些圣经不断宣告祂既是天主也是天使。此外，在这一切之上还加上：正如神圣圣经常宣告祂既是天使又是天主，同样神圣圣经也宣告祂既是人又是天主，藉此表达祂将要成为的，并当时以预象描绘了祂在实体真理中将要成为的。\Quote{因为，}经上说，\Quote{雅各伯独自一人留下；有一个人与他摔跤直到天亮。那人见自己不能胜过他，就摸了雅各伯大腿的阔处，他正与雅各伯摔跤，雅各伯也与那人摔跤，那人对他说：\InnerQuote{让我走，因为天亮了。}他说：\InnerQuote{你不祝福我，我就不让你走。}那人说：\InnerQuote{你叫什么名字？}他说：\InnerQuote{雅各伯。}那人对他说：\InnerQuote{你的名字不再叫雅各伯，而要叫以色列，因为你与天主与人角力，都得了胜。}}经上又加上说：\Quote{雅各伯给那地方起名叫天主的异象，因为他说：\InnerQuote{我面对面见了主，我的性命得以保全。}太阳升起照着他。他随后过了天主的异象，大腿却瘸了。}经上说，有一个人与雅各伯摔跤。如果这只是一个普通人，他是谁？他从哪里来？他为何与雅各伯争斗摔跤？有什么干预？发生了什么？如此大的争端和争斗的原因是什么？再者，为何雅各伯——他力量强大到能抱住与他摔跤的那个人，并向他求祝福——被断言因此请求，除非因为这场争斗是预象，预象基督与雅各伯子孙之间的争斗，而这事据说在福音中应验了？因为雅各伯的百姓与这个人争斗，在这场争斗中雅各伯的百姓显得更强，因为他们以不义战胜了基督：那时，由于他们犯下的罪行，犹豫退却，开始在自己信德和救恩的行走中严重瘸腿；虽然他们在定基督罪的事上显得更强，他们仍然需要祂的怜悯，仍然需要祂的祝福。但是，与雅各伯摔跤的那个人说：\Quote{此外，你的名字不再叫雅各伯，而要叫以色列。}如果以色列是看见天主的人，主美妙地显示出，当时与雅各伯摔跤的不仅是一个人，也是天主。的确，雅各伯看见了他与之摔跤的天主，尽管他当时在争斗中抱着那个人。为了不再有疑惑，祂亲自作出解释，说：\Quote{因为你与天主与人角力，都得了胜。}因此，同一位雅各伯，已经领悟奥秘的力量，并理解与他摔跤的那位的权威，就给他摔跤的地方起名叫天主的异象。他还在解释天主的异象时补充了理由：\Quote{因为，}他说，\Quote{我面对面见了天主，我的性命得以保全。}此外，他看见了天主，与他摔跤如同与一个人；但他确实作为征服者抱着那个人，虽然作为卑微者，他如同从天主那里求祝福。因此，他与天主和人摔跤；如此，那场争斗真实地预象了基督与雅各伯百姓之间的争斗，并在福音中应验：虽然百姓占上风，但因被显为有罪而显为低下。谁会犹豫承认基督——这摔跤的预象在他身上应验——不仅是一个人，也是天主，因为就连那摔跤的预象本身也似乎证明他是人和天主？然而，即使在此之后，同一位神圣圣经仍公正地不停止称天使为天主，并宣告天主为天使。因为当这雅各伯将要祝福若瑟的儿子默纳协和厄弗辣因时，他交叉双手按在孩子们的头上，说：\Quote{从我幼年直到今日牧养我的天主，救我脱离一切祸患的天使，祝福这两个少年。}甚至到如此程度，他确认那同一位是他曾称为天主的，也是天使，以致在讲论结束时，当他所说\Quote{祝福这两个少年}时，表达了他所讲的那一位是同一的。因为如果他意在将一位理解为天主，另一位为天使，他就会在祝福中使用复数指两个位格；但现在他定义了祝福中单数的一位，由此他意图使人们理解同一位是天主和天使。但祂不能被理解为天主父；而作为天主和天使，祂可以被理解为基督。雅各伯通过将手交叉放在少年们头上，也表明祂是这祝福的作者，仿佛他们的父亲是基督，并藉此手形表明受难的预象和未来形式。因此，凡不避讳称基督为天使的，也当不避讳称祂为天主，当他认识到在祝福这些少年时，通过\Emph{交叉的}双手的预象所暗示的受难圣事，祂本身同时作为天主和天使被呼求。\par

% structure: p0062 source=h2 target=section reason=relative-heading-rank; base=h2

\section{第二十章。}

\Emph{从圣经证明基督曾被称为天使。但圣经其他部分也表明祂是天主。}\par

但如果某个异端者固执地反对真理，坚持在所有例子中认为基督原本是天使，或坚持必须这样理解，那么他在这方面也必须被真理的力量征服。因为如果天上、地上和地下的一切都隶属于基督，就连天使本身和所有其他受造物，凡隶属于基督的，都被称为众神，那么基督也正确地被称作天主。如果任何隶属于基督的天使都可以被称作神，而这样称呼也未被视为亵渎，那么对基督、天主子本身而言，称祂为神当然更合适。因为如果一位天使隶属于基督而被尊为神，那么基督——众天使都隶属于祂——就更应且更一致地被称作神。因为这不适合本性：对较小的所允许的，不应拒绝给较大的。因此，如果一位天使低于基督，而天使尚且被称作神，那么基督按逻辑更应被称作神，因为祂被发现比一位天使伟大和优越，而且比所有天使都伟大和优越。而\BibleQuote{天主持立于众神的会中，在诸神中间辨别诸神}，并且基督在不同时间站在会堂中，那么基督作为天主站在会堂中——即在诸神之间审判，祂对他们说：\BibleQuote{你们偏袒人的情面要到几时？}这就是说，祂谴责会堂中的人不实行公正的审判。此外，如果那些被责备和谴责的人，即使因任何理由而获得此名而不算亵渎，即被称为众神，那么祂更应被尊为天主，祂不仅被记载曾作为天主站在众神的会堂中，而且还在同一读经的权威下被启示为在诸神之间辨别和审判。但\Emph{即使}那些\BibleQuote{像王子一样跌倒}的仍被称为众神，那么祂更应被称作天主，祂不仅不像王子一样跌倒，反而亲自战胜了邪恶的创始者和元首。世上有什么理由，尽管他们说这名字甚至给了梅瑟，因为经上说：\BibleQuote{我立你作法郎的天主}，却要把这名字拒绝给基督？基督被宣告为不仅是法郎\Emph{唯一}的主和天主，而且也是所有受造物的主和天主。在前一种情况下，这名字是有所保留地给予的，在后一种情况则慷慨地给予；在前者是有限度地，在后者则超越一切限度：\Quote{因为，}经上说：\BibleQuote{父不赐给子有限量，因为父爱子。}在前者是暂时的，在后者则不受时间限制；因为祂接受了神圣名字的能力，超越一切事物并永远有效。但如果一个人接受了另一个人的权力，在这有限权力下仍毫不迟疑地获得了神的名字，那么那位连梅瑟本身都有权柄的，岂不更应被相信获得了那名字的权威？\par

% structure: p0065 source=h2 target=section reason=relative-heading-rank; base=h2

\section{第二十一章。}

\Emph{同一天主尊威再次由其他圣经在基督内得到证实。}\par

诚然，我能用全部神圣圣经来阐述这主题，并且可以说，为了揭示基督神性的显现，我能够发动一场完美的\Emph{经文}森林。然而，我主要不是要针对这种特殊的异端形式进行辩论，而是要阐明关于基督位格的真理规条。不过，尽管我必须赶快处理其他事项，但我认为不可忽略这一点：在福音中，主为表明祂的尊威，曾宣告说：\BibleQuote{拆毁这座圣殿，三天内我要把它重建起来。}或者，在另一段经文、另一主题上，祂宣告：\BibleQuote{我有权舍弃我的生命，也有权再取回它；这是我从我父所接受的命令。}那么，是谁说祂能舍弃祂的生命，或能自己恢复祂的生命，因为祂从父接受了它？或者说，谁能再次复活并重建祂被毁的身体的圣殿，除非因为祂是来自父的圣言，与父同在，\BibleQuote{万物是借着祂造成的，凡受造的，没有一样不是由祂造成的；}祂是父工作的模仿者，是父权能的模仿者，\BibleQuote{是不可见的天主的肖像，}\BibleQuote{是自天降下的，}祂见证了自己所见所闻的事；祂\BibleQuote{来不是为执行自己的意愿，而是为执行那派遣祂来的父的旨意，}父正是为此目的派遣了祂，使祂成为\BibleQuote{伟大的参谋使者}，向我们揭示天上奥秘的律法；并且作为圣言成了血肉，居住在我们中间，这位基督被证明不单是人，因为祂是人之子，而且也是天主，因为祂是天主子？如果宗徒称基督为\BibleQuote{一切受造物的首生者}，祂怎能是一切受造物的首生者，除非按照祂的神性，圣言在一切受造物之前已由父而生？除非异端者如此接受，否则他们将不得不证明基督之人是一切受造物的首生者，这是他们做不到的。因此，要么祂在一切受造物之前，以便成为一切受造物的首生者，而祂不单是人，因为人在一切受造物之后；要么祂只是人，那么祂就在一切受造物之后。那么，祂又怎么是一切受造物的首生者呢？除非因为祂是那在一切受造物之前的圣言；因此，作为一切受造物的首生者，祂成了血肉，居住在我们中间，即取了那在一切受造物之后的人性，并与之同住，且在他内、在我们内，如此基督的人性不被除去，神性也不被否认。因为如果祂只在一切受造物之前，人性就从祂身上被除去；但如果祂只是人，那在一切受造物之前的神性就受到妨碍。因此，这两者在基督内联合、结合，彼此联系。并且正确地说，正如祂内有着超越受造物的东西，神性与人性的\Emph{结合}似乎在祂内得到了保证：为此，那被宣告为\BibleQuote{天主与人之间的中保}的，被启示为在祂自身内结合了天主和人。如果同一宗徒论及基督说：\BibleQuote{脱去了\Emph{肉身}，祂解除了权势，公然仗着祂自己凯旋了它们，}那么他一定不是毫无意义地提出脱去肉身，除非因为他希望人理解在复活时祂又穿上了肉身。那么，谁是那位这样脱去又穿上肉身的呢？让异端者去查考吧。因为我们知道，天主的圣言穿上了血肉的实体，然后又脱去了同一身体性的材料，这材料祂在复活中又取回并如同衣服一样穿上。然而，如果基督只是人，祂既不能脱去也不能穿上人性：因为人从未被剥夺或穿上自己。因为那由一个他者取去或穿上的，必须是某种别的东西，无论它是什么。因此，合理地，是天主圣言脱去了肉身，并在复活中又穿上了它，因为祂脱去它，是因为在出生时祂曾穿上了它。因此，在基督内，是天主穿上并也必须脱下，因为那穿上的也同样必须是那脱下的；而作为人，祂像是穿上一件由紧凑身体构成的外衣，又脱下它。因此，结果，正如我们所说，祂是天主圣言，被启示为时而穿上、时而脱下\Emph{肉身}。因为这一点，祂早在祝福中预言：\BibleQuote{祂将在酒中洗净祂的衣服，在葡萄的血中洗净祂的衣裳。}如果在基督内衣服是肉身，衣裳本身就是身体，那么就应问那以身体为衣裳、以肉体为衣服者是谁？因为对我们来说，明显的是肉体是衣服，身体是圣言的衣裳；祂洗净了祂身体的实体，在血中，即在酒中，借着祂所受的苦难，在祂所承担的人性内，净化了肉体的材料。因此，如果祂确实被洗净，祂就是人，因为被洗净的衣服是肉身；但洗净者是天主圣言，祂为了能洗净衣服，成了衣服的承受者。正确地说，从那被取来以便洗净的实体，祂被揭示为人，正如从洗净衣服的圣言的权威，祂被显示为天主。\par

% structure: p0068 source=h2 target=section reason=relative-heading-rank; base=h2

\section{第二十二章}

\Emph{基督内具有同一神性，祂再次以其他圣经断言此事}\par

\BibleQuote{祂虽具有天主的形体，却没有将自已与天主同等视为僭越；反而空虚自己，取了奴仆的形体，成了人的模样；既以人的形态出现，就贬抑自己，听命至死，且死在十字架上。因此，天主极其举扬祂，赐给祂一个超越众名的名字，为使因耶稣的名字，天上、地上和地下的一切，无不屈膝叩拜；并且众口皆宣认耶稣基督是主，以光荣天主圣父。} 他说：\Quote{祂虽具有天主的形体，} 如果基督仅是凡人，祂就会被说成具有天主的\Quote{形像}，而非\Quote{形体}。因为我们知道，人是按天主的形像或相似而造的，并非按天主的形体。那么，我们所说的那位具有天主形体的天使是谁呢？但在天使身上，我们也读不到天主的形体，唯独这一位，祂是超越一切的首领与君王——天主子，天主的圣言，祂效法祂圣父的一切作为，因为祂自己行事正如祂的父。祂——正如我们所宣称的——具有天主圣父的形体。祂合理地被称为具有天主的形体，因为祂本身超越万有，拥有治理一切受造物的神性权能，并效法圣父而成为天主。然而，祂从自己的圣父那里获得这一点：祂应为万有的天主，应为主，并被生为出自祂本身的天主，具有天主圣父的形体，而为人所知。因此，祂虽具有天主的形体，却未将自已与天主同等视为僭越。因为尽管祂记得自己是出自天主圣父的天主，却从不将自己与天主圣父相比较或等同，祂深知自己来自圣父，并且祂所拥有的那一切，皆因圣父赐予了祂。由此，最终，无论在降生成人之前，还是在取了肉身之后，甚至在复活之后，祂都向圣父表示完全服从，并一如既往地服从。由此证明，祂认为某种神性\Emph{的宣称}乃是僭越，即与天主圣父平等；但另一方面，祂顺服并服从于圣父的一切统治与旨意，甚至满足于取了奴仆的形体——即成为人；而血肉与身体的实质，按祂的人性而言，来自祂先祖罪过的束缚，祂通过降生而承担，那时祂更是空虚了自己，因为祂没有拒绝承受人性固有的脆弱。因为如果祂仅是作为人而生，在这方面祂就不会空虚；因为人出生时是增加，而非空虚。因为正如我们所说，在开始成为祂之前所不能拥有的事物时，祂并非空虚，反而是增加和充实。但如果基督在降生、取了奴仆的形体时是空虚的，祂又怎能仅仅是凡人呢？关于祂，在祂降生时，更能真实地说祂是充实而非空虚，除非是因为圣言的权能暂时安息在取了人性之中，没有发挥其真实的力量，它自我贬抑，暂时放下自己，承担起所取的人性？它在降到伤害与凌辱、忍受憎恶、经历不堪之事时空虚自己；然而这谦卑立刻带来了崇高的赏报。因为祂\Quote{获得了超越众名的名字}，我们确知这无非是天主的名字。因为唯独天主超越万有，因此那属于超越万有之上者的名字超越众名；这个名自然属于那位，祂虽具有\Quote{天主的形体，却未将自已与天主同等视为僭越}。因为如果基督不是天主，那么天上、地上和地下的一切就不会在祂的名下屈膝；\BibleQuote{天上、地上和地下的一切} 可见与不可见的，甚至一切受造之物，就不会臣服或置于人之下，既然它们可能记得自己先于人。因此，既然基督被称为具有天主的形体，并显示祂按肉身受生而空虚了自己；既然宣告祂从圣父领受了那超越众名的名字；既然显示在祂的名下，\BibleQuote{天上、地上和地下的一切，无不屈膝叩拜}；并且此事本身被宣称是为了光荣天主圣父；那么，祂就不仅仅是人，因为祂服从圣父至死，且死在十字架上；而且，从这些更高之事对基督神性的宣告中，基督耶稣被显示为主和天主，这是异端者所不愿承认的。\par

% structure: p0071 source=h2 target=section reason=relative-heading-rank; base=h2

\section{第二十三章}

\Emph{这一点如此明显，以至于有些异端者认为祂是天主圣父，另一些则认为祂只是没有血肉的天主。}\par

在此，我也许可以许可也从其他异端者方面收集论据。这是一种实质性的证明，甚至从敌人那里收集而来，以便甚至从真理的敌人本身来证明真理。因为祂在圣经中被宣告为天主，这一点是如此明显，以致许多异端者，被这神性的伟大和真理所感动，过度夸大对祂的尊崇，竟敢宣告或认为祂不是子，而是天主圣父本身。虽然这与圣经的真理相悖，但它仍然是基督神性的一个伟大而卓越的论据：祂是如此天主，除作为天主之子、由天主所生之外，以至于许多异端者——正如我们所说——如此接受祂为天主，以致认为祂必须被宣告为父，而不是子。因此，让我们考虑祂是否是天主\Emph{或不是}，因为祂的权威如此影响了一些人，以至于如我们上面所说，他们以为祂就是天主圣父本身，并以如此猛烈和洋溢的方式承认了基督内的神性——被基督内明显的神性所迫——以至于他们认为他们读到的那位为子的，因察觉到祂是天主，就必定是父。此外，其他异端者如此接受了基督明显的天主性，以至于说祂没有肉体，并撤去祂所取的全部人性，以免——如他们所想——因将人的出生与祂相联而削弱祂天主之名的能力。然而，我们并不赞同这一点；但我们引用它作为论据来证明基督是天主，达到如此程度：一些人除去人性，认为祂只是天主；另一些人认为祂就是天主圣父本身；而圣经的理性和比例显示基督是天主，但作为天主子；并且祂一直是人子，曾被天主所收纳，因此也必须相信祂也是人。因为如果祂来到人间，成为天主与人之\OriginalTerm{中保}{Mediator}，祂就必须与人同在，圣言必须成为血肉，以便在祂自身内将天上与地下之事的和好联结起来，通过在自己内联合两性的担保，将天主与人和人与天主结合；这样，天主子因取肉身而成为人子，人子因接受天主的圣言而成为天主子，这是合理的。这极其深奥隐秘的奥秘，在万世之前为人类救恩所预定，在耶稣基督——既是天主又是人——身上得以实现，以使人类能够获得永生救恩的享乐。\par

% structure: p0074 source=h2 target=section reason=relative-heading-rank; base=h2

\section{第二十四章}

\Emph{因此，那些人因认为天主子与\OriginalTerm{人子}{Son of Man}之间没有区别而错误，因为他们\Emph{错误地}理解了圣经。}\par

但是，那种异端错误的材料，依我判断，源于此：他们认为天主子和人子之间没有区别；因为如果加以区分，耶稣基督就很容易被证明既是人又是天主。因为他们主张，同一个人、人子，也显为天主子；人、肉体和那同一脆弱的实体，也可被称为天主自己之子。因此，既然在人子和天主子之间看不出区别，反而人子本身被断言为天主子，那么同一基督和天主子就被断言仅仅是人是；由此他们力图排除\BibleQuote{圣言成了血肉，寄居在我们中间}。并且你要称祂的名字为厄玛奴耳；翻译出来就是：天主与我们同在。因为他们提出并引出路加福音中所说的，由此他们力图坚持的不是真理，而只是他们想要的东西：\BibleQuote{圣神要临于你，至高者的能力要庇荫你，因此那要诞生的圣者，将称为天主子}。那么，他们说，如果天主的天使对玛利亚说\BibleQuote{那要诞生的圣者}，血肉和身体的实体是来自玛利亚；但他表明这实体，即那从她诞生的圣者，是天主子。他们说，人本身，那身体的肉；那被称为圣者的，本身就是天主子。那也当圣经说\BibleQuote{圣者}时，我们应当理解为基督是人，人子；而当它向我们呈现天主子时，我们应当感知，不是人，而是天主。然而，神圣的圣经容易揭露和揭穿异端者的欺诈和诡计。因为如果只是这样说：\BibleQuote{圣神要临于你，至高者的能力要庇荫你；因此那从你诞生的圣者，将称为天主子}，也许我们不得不用另一种方式与他们争辩，并寻求其他论据，拿起其他武器来战胜他们的陷阱和诡计；但由于圣经本身，充满天上的丰盈，解除了这些异端者的诽谤，我们容易依靠所写的内容，毫不犹豫地克服那些错误。因为它没有像我们已经陈述的那样说：\BibleQuote{因此那从你诞生的圣者}；而是加上了连词，因为它说：\BibleQuote{因此那从你诞生的圣者，\Emph{也}}，这样，就清楚表明那从她诞生的圣者——即那血肉和身体的实体——首先不是天主子，而是其次、在第二位；但首先，天主子是天主的圣言，由那位天使所说的圣神降生成人。因为祂是出自天主本身的合法天主子；而祂，在摄取那圣者、将人子联结到自己、将祂拉向自己并转移到自己身上时，通过祂的连接和混合的联合，负责并使祂成为天主子，而祂本性并不是，以致天主子这名号的原始原因是在于降临而来的主之神那里，而这人子那里的名号只是延续；因此，合理地，祂成了天主子，虽然原本祂不是天主子。因此天使，看到那安排，并顾及那奥秘的秩序，并没有混淆一切以致不留区别的痕迹，而是通过说\BibleQuote{因此那从你诞生的圣者，也将称为天主子}确立了区别；免得，如果他不用他的天平安排好那划分，而是让事情全部混乱地混合，那真会给异端者提供机会，宣称人子，就祂是人而言，与天主子和人是一样的。但现在，他分别解释了如此伟大奥秘的规约和理由，他在说\BibleQuote{那从你诞生的圣者，将称为天主子}时清楚地表明：天主子降临了，并且祂在将人子接入自己内时，从而使人子成为天主子，因为天主子将祂与自己联合并连接。这样，人子在其生诞中依附于天主子，通过那混合持有祂本性所不能拥有的作为抵押和派生。因此，通过天使的话，区别被确立，反对异端者的愿望，在天主子与人之间；然而又结合他们，迫使他们理解基督人子是人，也接受天主子与人天主子；即天主的圣言，如经上所记，是天主；从而承认主耶稣基督，可以说是两边联结，两边编织并生长在一起，通过相互联盟的绑定，在两种实体的同一协议中联合——人与天主，凭借宣告这同一事实的圣经的真理。\par

% structure: p0077 source=h2 target=section reason=relative-heading-rank; base=h2

\section{第二十五章。}

\Emph{然而，并不由此推论，因为基督死了，就必须也接受天主死了；因为圣经表明基督不仅是天主，也是人。}\par

因此，他们说，如果基督不但是人，也是天主——而圣经告诉我们，祂为我们死了，且复活了——那么圣经就教导我们相信天主死了；或者，如果天主不死，而基督被说成死了，那么基督就不是天主，因为不能承认天主死了。如果他们能够理解或曾经理解他们所读的，就绝不会以如此危险的方式说话。但错误的愚昧总是急于下滑，那些背弃了合法信仰的人甚至下滑到危险的结局，这并不是什么新鲜事。因为如果圣经只阐述基督是天主，而祂的本性中并没有搀和人性软弱，那么他们这个复杂的论点或许有些道理。如果基督是天主，而基督死了，那么天主就死了。但当圣经断定（我们已多次指出），祂不但是天主，也是人，那么便得出结论：那不朽的可以视为保持未受损。因为谁不能明白天主性是不受苦的，虽然人性软弱是易于受苦的？因此，当基督被理解为既是天主又是人——因为\BibleQuote{圣言成了血肉，住在我们中间}——谁能不靠任何老师和解释者，轻易自己领会到：在基督内死的并非那属于天主的，而是那属于人的？因为，如果在基督内的天主性不死，而只有肉身的实体被毁灭，那么在其他人身上——他们不单是肉身，而是肉身与灵魂——的确只有肉身遭受消耗和死亡的侵袭，而灵魂却显示为不受败坏，超越毁灭和死亡的律法，又当如何呢？因为我们的主自己也说过，劝勉我们殉道并轻视一切人间权力：\BibleQuote{你们不要害怕那些杀害肉身，而不能杀害灵魂的}。但如果不朽的灵魂在任何其他人身上都不能被杀害或杀死，尽管身体和肉身本身能被杀死，那么，天主的\OriginalTerm{圣言}{Logos}和天主在基督内，虽然肉身和身体被杀死，岂不更是绝对不能被处死！因为如果在任何一个人身上，灵魂具有不朽的卓越性，以致不能被杀害，那么天主的圣言的尊贵更拥有不被杀害的权能。因为如果人的权能不能杀害天主的圣能，如果人的残忍不能毁灭灵魂，那么它更不能杀害天主的圣言。因为就如灵魂本身，它是由天主的圣言所造的，不被人类杀害，那么当然更会相信，天主的圣言不能被毁灭。而如果人类血腥的残忍对人所能做的，不过是杀害身体，那么它必定更没有能力对付基督，超乎杀害身体的方式之外。因此，既然从这些考量可以推断出，在基督内只有人性被处死，那么祂内的圣言似乎并没有被拖入死亡。因为如果亚巴郎、依撒格和雅各伯——他们被公认只是人——显明是活着的——因为祂说，\BibleQuote{他们都生活于天主；}死亡在他们身上并没有毁灭灵魂，虽然它分解了身体本身：因为死亡能在身体上施展其能力，却不能对灵魂施展；因为他们身上的那部分是可朽的，因此死了；另一部分是不朽的，因此被认为没有被熄灭；因此他们被肯定地说生活于天主——那么死亡在基督身上更能只对祂身体的质料施展能力，而对圣言的天主性，却不能施展自身。因为当不朽的权能介入时，死亡的能力就被粉碎了。\par

% structure: p0080 source=h2 target=section reason=relative-heading-rank; base=h2

\section{第二十六章}

\Emph{此外，他反对\OriginalTerm{撒贝尔异端}{Sabelliani}，证明圣父是一位，圣子是另一位。}\par

但借着基督由圣经的神圣权威被证明不但是人，也是天主这一机会，其他异端者兴起，设法削弱在基督内的信仰立场；他们想借此表明基督就是天主圣父，因为祂被断言不但是人，也被宣告是天主。因为他们这样说：如果天主是唯一的，而基督是天主，那么他们说：如果圣父和基督是同一天主，基督就会被称为圣父。在此他们被证明是错误的，不认识基督，而是跟随名字的声音；因为他们不愿祂是圣父之后的第二位，而是圣父本身。既然这些事容易回答，就用几句话来说。因为谁不承认圣子的位格是圣父之后的第二位，当他读到圣父对圣子说：\BibleQuote{让我们照我们的肖像，按我们的模样造人；}并且之后记载：\BibleQuote{天主造了人，按天主的肖像造了他？}或者当他手中持有：\BibleQuote{上主从天上从主那里降下火与硫磺，落在索多玛和哈摩辣？}或者当他读到（作为已说的）对基督说：\BibleQuote{你是我的儿子，我今日生了你。你向我请求，我必将万民赐你为产业，地极为你的基业？}或者当那位可爱的作者也说：\BibleQuote{上主对我主说：你坐在我右边，直到我使你的仇敌作你的脚凳？}或者当他展开依撒意亚的预言，发现这样写着：\BibleQuote{上主对我主基督这样说？}或者当他读到：\BibleQuote{我不是从天上降下来为行我的旨意，而是行那派我来者的旨意？}或者当他发现写着：\BibleQuote{因为那派我来的，比我更大？}或者当他思考这段：\BibleQuote{我往我父和你们父那里去，往我的天主和你们的天主那里去？}或者当他发现与其他的并列：\BibleQuote{并且，在你们的法律中记载着两个人的见证是真的。我为自己作证，派我来的父也为我作证？}或者当天上的声音是：\BibleQuote{我已经光荣了祂，我还要再光荣祂？}或者当伯多禄回答说：\BibleQuote{你是永生天主的儿子？}或者当主亲自认可这启示的圣事，说：\BibleQuote{有福的是你，西满巴尔约纳，因为不是血肉启示了你，而是我在天之父？}或者当基督自己表达：\BibleQuote{父啊，在未有世界之前，我同你所有的光荣，求你使我与你同享？}或者当同一主说：\BibleQuote{父啊，我知道你常俯听我；但我因周围站着的群众说了这话，使他们信是你派了我来？}或者当准则的定义由基督自己确定，说：\BibleQuote{永生就是：认识你，唯一的真天主，和你派来的耶稣基督。我已在地上光荣了你，完成了你派我做的工？}或者当此外由同一主断言并说：\BibleQuote{一切由我父交付给我？}或者当坐在父的右边被宗徒和先知证明？若我努力收集所有与此相关的段落，我将会忙碌不已；因为圣经，不仅是旧约，也是新约，处处显示祂由父所生，万有藉祂而造，没有祂什么也不能造，祂总是顺服且顺服父；祂永远掌管万有，但这是作为交付、赐予、由父自己许可给祂的。有什么能如此明显的\Emph{证据}证明这不是父，而是子？正如祂被展示为顺服于天主父，除非，若祂被认为是父，则基督可能被说成是顺服于另一位天主父？\par

% structure: p0083 source=h2 target=section reason=relative-heading-rank; base=h2

\section{第27章}

\Emph{他巧妙地回应了异端者用来辩护自己观点的一段经文。}\par

然而，因为他们常常用这段经文来攻击我们，即经上记载的：\BibleQuote{我与父原是一体}，我们在这一点上也能轻易驳倒他们。因为如果如同异端者所想的，基督就是圣父，祂就该说\BibleQuote{我与父原是一体}。但当祂说\Quote{我}，接着又提到父说\Quote{我与父}时，祂就切断并区分了祂自己（即圣子）的\OriginalTerm{位格}{persona}的特性与圣父的权威，不仅是在名字的发音上，而且是在权能分配的秩序上，因为如果祂认为自身就是父，祂大可以说\Quote{我是父}。由于祂说的是\Quote{一}物，让异端者明白，祂说的不是一位位格。因为放在中性里的\Quote{一}，暗示的是社会性的和谐，而非位格的同一。祂被称为中性的\OriginalTerm{一（中性）}{unum (neutrum)}，而非阳性的\OriginalTerm{一（阳性）}{unus (masculinum)}，因为这句话不是指数量，而是就与另一位的关联而言的。最后，祂又加上\Quote{我们是}，而不是\Quote{我是}，以此表明，祂说\Quote{我与父}乃是两个位格。此外，祂说\Quote{一}，指的是合一、判断的一致以及爱的关系本身，因此圣父与圣子在意志、爱和情感上是一体；而且因为祂来自圣父，无论祂是什么，祂都是圣子；然而区别仍然存在：圣子不是父，因为圣父不是子。因为如果祂认为独一的圣父成了圣子，祂就不会加上\Quote{我们是}。此外，宗徒\Person{保禄}也理解这种合一的合一，同时保留位格的区分：因为他在致格林多人书函中写道：\BibleQuote{我栽种了，阿波罗浇灌了，然而使之生长的，却是天主。所以，栽种的不算什么，浇灌的也不算什么，只在乎那使之生长的天主。栽种的和浇灌的原是一体。}谁不明白\Person{阿波罗}是一个位格，\Person{保禄}是另一个位格，而\Person{阿波罗}和\Person{保禄}不是同一位格？而且，他们各自提到的职务也不同：一位是栽种的，另一位是浇灌的。但\Person{保禄}宗徒提出这两个人，不是作为一个位格，而是作为一体；所以，尽管\Person{阿波罗}是一个，\Person{保禄}是另一个，就位格的区分而言，但就他们的合一而言，两者都是\Quote{一}。因为当两个人有一个判断、一个真理、一个信德、一种同样的宗教、同样的敬畏天主时，他们虽然仍是两个位格，却是一体；他们在拥有同一个心志上是相同的。既然位格之考量将彼此分开的，宗教之考量又将他们合为一。虽然他们实际上不是同一个人，但在感受相同上，他们是相同的；虽然他们是两个人，仍然是一体，如同在信德中有共融，尽管位格上有差异。此外，当主的话激起犹太人的无知，他们急忙拿起石头，说：\BibleQuote{我们为了善事并不用石头砸你，而是为了亵渎的话，并且因为你是人，却把自己当作天主。}主便以祂称自己为天主或愿意人如此理解的原则来确立区分，并说：\BibleQuote{那么，父所祝圣并派遣到世界上来的，你们就说：你说亵渎的话，因为我曾说过：我是天主子吗？}这里祂也提到祂有父。因此，祂是子，不是父；因为如果祂认为自己是父，祂就会承认自己是父；而祂宣告自己是由父所祝圣的。那么，既然从父领受祝圣，祂就比父为小。因此，比父小的，就不是父，而是子；因为如果祂是父，祂就会给予祝圣，而非领受。但现在，祂藉着宣称从父领受祝圣，藉着证明自己比父小这一事实，表明祂是子，而不是父。此外，祂说祂被派遣；这样，主\Person{基督}以服从被派遣而来，证明祂不是父，而是子；如果祂是父，祂肯定会派遣，但祂被派遣，所以祂不是父，免得父在被派遣的情况下成为受制于另一位天主。此后，祂又加上一句可以消解一切模糊并扑灭错误争议的话：在祂讲论的末段，祂说：\BibleQuote{你们说：你说亵渎的话，因为我曾说过：我是天主子。}因此，如果祂明明确证自己是天主子，而不是父，那么，与主\Person{基督}自己的见证相反，掀起神性与宗教的争议，声称\Person{基督耶稣}就是父，就是极大的鲁莽和疯狂，因为祂已被证明是子，而不是父。\par

% structure: p0086 source=h2 target=section reason=relative-heading-rank; base=h2

\section{第二十八章}

\Emph{他也证明对\Person{斐理伯}所说的话对\OriginalTerm{撒伯里乌派}{Sabellians}毫无用处。}\par

此外，我也要加上那一种观点，即异端者虽因失去某种看见特殊真理与光明的能力而沾沾自喜，却承认自己错误的完全盲目。因为他屡次、频繁地提出反对说，经上记载：\BibleQuote{斐理伯！这么长久的时候，我和你们在一起，而你们还不认识我吗？谁看见了我，就是看见了父。} 但让他学习他所不明白的事。斐理伯受到责备，是正确且应当的，因为他曾说：\BibleQuote{主！把父显示给我们，我们就心满意足了。} 因为他何时曾从基督那里听说过，或学到过基督就是父呢？虽然，另一方面，他经常听到且学到的是：祂是子，而不是父。因为主说：\BibleQuote{你们若认识我，也就认识我的父；从今以后，你们认识祂，并且已经看见祂。} 祂这样说，并非愿意自己被理解为父，而是\Emph{暗示}：凡是完全、圆满、以全部信德和全部虔诚亲近天主圣子的人，必定会藉着祂所相信的圣子本身，到达父那里，并看见祂。因为祂说：\BibleQuote{除非经过我，谁也不能到父那里去。} 因此，他不仅会到天主父那里，认识父本身；而且，他还应当这样持守，在心里如此认定，即他已不仅认识，而且看见了父。因为神圣的圣经往往将尚未成就的事说成已经成就，因为将来必会如此；而那些必将发生的事，它不预言为将来，却叙述为已成就。这样，虽然在依撒意亚先知时代基督尚未诞生，他却说：\BibleQuote{因为有一个婴孩为我们诞生了；} 虽然玛利亚尚未受孕，他却说：\BibleQuote{我便走近了女先知；她怀孕，生了一个儿子。} 当基督尚未将父的意愿显明时，经上说：\BibleQuote{祂的名要称为\OriginalTerm{伟大的谋士天使}{Angel of Great Counsel}。} 当祂尚未受难时，祂宣告：\BibleQuote{祂如同一只被牵去宰杀的羊。} 虽然十字架尚未存在，祂说：\BibleQuote{我整天向悖逆的百姓伸开双手。} 虽然尚未被人戏弄地给祂喝，圣经却说：\BibleQuote{我渴了，他们竟拿醋给我喝。} 虽然祂尚未被剥去衣服，祂说：\BibleQuote{他们为我的衣服拈阄，数算我的骨头，扎了我的手和脚。} 因为神圣的圣经有先见之明，将必然成就之事说成已成就，将将来无疑要发生的事当作已完成的叙述。因此，主在当下经文中说：\BibleQuote{从今以后，你们认识祂，并且已经看见祂。} 祂说凡跟随圣子的人必看见父，并非说圣子本身是那被看见的父，而是说凡愿意跟随祂、作祂门徒的人，必能得到看见父的酬报。因为祂也是天主父的肖像；因此，此外还加上：\BibleQuote{父怎样工作，子也照样工作。} 子效法父的一切作为，以致当人看见祂在一切作为中效法那不可见的父时，便可视为如同看见了父。但如果基督自己就是父，祂为何立即接着说：\BibleQuote{凡信我的，我所做的工程，他也要做，并且还要做比这些更大的工程，因为我往父那里去？} 祂又加上：\BibleQuote{如果你们爱我，就要遵守我的命令；我要求父，祂必赐给你们另一位\OriginalTerm{护慰者}{Comforter}。} 此后祂又说：\BibleQuote{谁爱我，必遵守我的话，我父也必爱他，我们要到他那里去，并要在他那里作我们的住所。} 此外，祂还加上：\BibleQuote{但那\OriginalTerm{护慰者}{Advocate}，就是父因我的名所要派遣的圣神，祂必要教训你们一切，也要使你们想起我对你们所说的一切。} 祂又说出那段话，表明自己是子，并合理地接着说：\BibleQuote{你们若爱我，就该喜欢我往父那里去，因为父比我大。} 但当祂继续这样说时，我们\Emph{该说什么}？\BibleQuote{我是真葡萄树，我父是园丁。凡在我身上不结实的枝条，祂便剪掉；凡结实的，祂就清理，使之结更多的果实？} 祂仍然坚持，并加上：\BibleQuote{正如父爱了我，同样我也爱了你们；你们应存留在我的爱内。如果你们遵守我的命令，你们便存留在我的爱内，正如我遵守了我父的命令而存留在祂的爱内一样。} 此外，祂还说：\BibleQuote{我却称你们为朋友，因为凡由我父听来的一切，我都显示给你们了。} 而且，祂还加上这一切：\BibleQuote{但他们要因我的名字对你们做这一切，因为他们不认识那派遣我来的。} 那么，这些在先前的话语之后，显然证明祂不是父而是子，主绝不会加上这些话——如果祂心里认为祂就是父，或愿意自己被理解为父——除非祂要宣布这一点：即每一个人在藉着子看见天主父的肖像时，从此应当认为如同看见了父；因为凡相信子的人，可操练注视这相似，从而习惯于在相似中看见天主性，并前进，直至达到对全能天主父的圆满注视。而谁将这真理吸收进思想和灵魂，并相信一切必如此成就，他现在就已如同看见了那\Emph{将来}要看见的父；他可以这样认为，如同他确实拥有那他知道总有一天会拥有的一切。但如果基督自己就是父，为什么祂将已经赐予和给予的酬报许诺为将来呢？因为祂说：\BibleQuote{心里洁净的人是有福的，因为他们要看见天主。} 这话被理解为许诺对父的默观与看见；因此，祂当时尚未赐予这个；因为如果祂已经赐予，为什么还要许诺？如果祂是父，祂就已经赐予了：因为祂被看见，也被触摸。但是，当基督自己被看见被触摸时，祂仍然许诺，并说心里洁净的人要看见天主，祂正是藉此证明，那当时在场者并不是父，因为祂被看见，却许诺心里洁净的人将看见父。因此，许诺这事的不是父，而是子；因为是子的许诺那尚待看见的事；如果祂是父，祂的许诺就是多余的。因为为什么祂向心里洁净的人许诺他们将要看见父，如果当时在场的人已经看见基督当作父呢？但由于祂是子，不是父，所以祂被看见时，正确地说，是作为子被看见，因为祂是天主的肖像；而父，因为是不可见的，被许诺和指出将由心里洁净的人看见。那么，针对那异端者，提出这些论点就足够了；关于许多事，只是简略地说了几句。因为如果我们想要更充分地讨论那异端者，就会展开一片既广阔又宽广的田地；因为他在这两点上，如同被挖去了双眼，完全被他教义的盲目所压倒。\par

% structure: p0089 source=h2 target=section reason=relative-heading-rank; base=h2

\section{第二十九章}

\Emph{祂接下来教导我们，信仰的权威命令我们在相信圣父和圣子之后，也要相信圣神，并列举圣神在圣经中的运行。}\par

此外，理智的秩序和在言语的措置及主的圣经中信德的权威，在这些事之后也劝诫我们相信那曾应许给教会、并在指定的时期赐下的圣神。因为祂曾由岳厄尔先知应许，却由基督赐下。先知说：\BibleQuote{在末日，我要将我的神倾注在我的仆人和婢女身上。} 主说：\BibleQuote{你们领受圣神吧！你们赦免谁的罪，谁的罪就得赦免；你们保留谁的罪，谁的罪就被保留。} 但主基督有时称这圣神为\Quote{护慰者}，有时又称祂为\Quote{真理之神}。祂并非在福音中是新的，也非新近才赐下；因为正是祂在先知中指责百姓，在宗徒中赐予他们向外邦人的呼吁。因为前者应受指责，因为他们藐视法律；而外邦人中信者应受圣神的辩护之助，因为他们热切渴望达到福音的法律。的确，在圣神内有不同种类的职务，因为在时代中有不同的时机的次序；然而，因此，执行这些职务的祂并非不同，也非在如此行事时是另一位，而是同一、不变的，按照时代、时机和事物的冲动分施祂的职务。此外，保禄宗徒说：\BibleQuote{具有同一圣神，正如经上所载：我信了，因此我说了话；我们也信，因此也说话。} 所以，那在先知和宗徒中的是同一的圣神，区别在于祂在先知的临在是暂时的，在宗徒内却是常存的；在先知的临在并非常在他们内，在宗徒内却是常在他们内；在先知的临在是有保留的分施，在宗徒内却是完全倾注；在先知的临在是吝啬地赐予，在宗徒内却是慷慨地赐予；在主复活之前尚未显明，在复活之后才赐予。因为祂说：\BibleQuote{我要请求父，祂必会赐给你们另一位护慰者，使祂永远与你们同在，就是真理之神。} 又说：\BibleQuote{当那护慰者，就是我从父那里要给你们派遣的，那发于父的真理之神来到时。} 又说：\BibleQuote{我若不去，那护慰者就不会到你们这里来；我若去了，我就派遣祂到你们这里来。} 又说：\BibleQuote{当那真理之神来到时，祂要引导你们进入一切真理。} 因为主将要升天，祂出于必要赐给门徒护慰者，以免他们成为孤儿（这是极不希望的），并舍弃他们没有护慰者和某种保护者。因为正是祂坚固了他们的心和思想，标明了福音的圣事，在他们内是神圣事物的启迪者；他们因被坚固，为了主的名不惧怕地牢或锁链，甚至践踏世界上的权势及其酷刑，因为他们从此被同一圣神武装和坚固，自身拥有这同一圣神所分赐的恩赐，并将它们归于教会、基督的净配，作为她的装饰。正是祂将先知安置在教会中，教导教师，引导言语，赐予能力和医治，行奇迹，常常分辨神类，提供治理的能力，建议计谋，并安排和布置任何其他\Emph{\OriginalTerm{神恩}{charismata}}的恩赐；从而在各处并在一切事上使主的教会得以完美和成全。正是祂，在我主受洗时，以鸽子的形像降来并停留在祂身上，住在基督内，圆满而完全，没有任何程度或部分的残缺；而是以其全部丰沛地分赐并派遣，使别人从祂那里得到祂恩宠的某种享受：整个圣神的源泉留在基督内，使人能从祂那里汲取恩赐和工作的溪流，而圣神丰盛地住在基督内。因为依撒意亚预言此事时，曾说：\BibleQuote{智慧和聪敏之神，要住在祂身上，超见和刚毅之神，知识和孝爱之神；敬畏主之神要充满祂。} 同样的事，他在另一处以主自己的名义说：\BibleQuote{主的神在我身上，因为祂给我傅了油，派遣我向贫穷人传报福音。} 达味同样说：\BibleQuote{因此，天主，你的天主，以喜乐之油傅了你，胜过你的同伴。} 保禄宗徒论祂说：\BibleQuote{谁若没有基督的圣神，谁就不属于基督。} \BibleQuote{主的神在那里，那里就有自由。} 正是祂与水分产生第二次出生，如同神圣繁衍的某种种子和天上诞生的祝圣，应许产业的凭据，以及永恒救恩的某种手印；祂能使我们成为天主的宫殿，并使我们适合于祂的家；祂以不可言喻的叹息为我们祈求天主的垂听，充满护慰的职务，显明我们辩护的职责——为我们的身体所赐的居住者和圣化者。祂在我们内为了永恒而工作，也能在不朽的复活中产生我们的身体，使它们习惯于与祂自己结合，具有天上的能力，并与圣神的神性永恒联合。因为我们的身体在祂内并藉着祂被训练迈向不朽，学习按照祂的法令以节制治理自己。因为正是祂\BibleQuote{与肉身相反，}因为\BibleQuote{肉身与圣神相抗。}正是祂抑制贪得无厌的欲望，控制无节制的情欲，熄灭非法的欲火，战胜鲁莽的冲动，驱退醉酒，制止贪财，驱逐奢华的宴乐，连系爱德，联结情谊，制止分裂，整顿真理的规则，克服异端者，赶出恶人，守护福音。关于祂，同一宗徒说：\BibleQuote{我们不是接受了世界的精神，而是那出于天主的圣神。}论祂，他欢跃地说：\BibleQuote{我也以为我有天主的圣神。}论祂，他说：\BibleQuote{先知的神是顺从先知的。}论祂，他也告诉我们：\BibleQuote{圣神明明地说，在末期有些人要背弃信德，听从诱惑的神、魔鬼的教训，这些人以虚伪说谎，良心被烙焦。}在这圣神内坚振的，\BibleQuote{没有人称耶稣为诅咒；}没有人曾否认基督是天主子，或拒绝天主造物主；没有人说出任何与圣经相反的私语；没有人制定其他亵渎的法令；没有人订立不同的法律。凡亵渎祂的，\BibleQuote{不得赦免，不但在今世，而且在来世也不得赦免。}正是祂在宗徒内为基督作证；在殉道者内显出他们宗教的恒久信实；在贞女内抑制他们封印贞洁的可钦之节操；在其他人内，守护主之教义的法律，使之不朽无玷；摧毁异端者，纠正乖戾者，定罪不信者，揭露伪装者；此外，责斥恶人，保守教会无腐无犯，在永恒童贞与真理的圣洁中。\par

% structure: p0092 source=h2 target=section reason=relative-heading-rank; base=h2

\section{第三十章}

\Emph{最后，尽管上述异端者从圣经所记载的内容中收集了他们错误的起源：虽然我们称基督为天主，称圣父为天主，但圣经并没有提出两位神，正如也没有提出两位主或两位教师一样。}\par

如今，关于圣父、圣子及圣神，这样简略地说这些，并扼要地陈明这些要点，而不以冗长的论述展开，便已足够。因为它们本可以更广泛地呈现并展开更详尽的辩论，因为全部旧约和新约都可引作见证，证实真信德如此成立。但是，由于异端者常与真理争辩，惯于延长对纯正传承及大公信德的争论，并因基督而发怒；此外，圣经也断言祂是天主，我们也如此相信；我们必须正当地——为使一切异端的诽谤从我们的信德中除去——就基督也是天主这一事实进行辩论，但这种辩论的方式不应损害圣经的真理，也不应损害我们的信德，即圣经如何宣示只有一位天主，以及我们如何持守和相信这一点。因为那些说耶稣基督自己就是天主圣父的人，以及那些认为祂只是人的人，都从那里收集了他们错误和邪恶的根源与理由；因为他们看到经上记载说\BibleQuote{天主是唯一的}，就认为他们只能通过假定要么基督只是人，要么祂真的是天主圣父，才能持守这样的观点。他们惯于如此连结他们的诡辩，以图证明自己的错误正当。因此，那些说耶稣基督是圣父的人如此辩论：如果天主是唯一的，而基督是天主，那么基督就是圣父，因为天主是唯一的。如果基督不是圣父，因为基督是天主圣子，那就似乎引入了两位天主，与圣经相悖。而那些主张基督只是人的人，则相反地得出结论：如果圣父是一位，圣子是另一位，并且圣父是天主，基督是天主，那么就不是一位天主，而是立刻引入了两位天主，即圣父和圣子；而如果天主是唯一的，那么基督必然是一个人，这样圣父才能正确地是唯一的天主。主确实仿佛被钉在两个盗贼之间，正如祂从前被安置的那样；因此祂从两边受到这些异端者的亵渎辱骂。但是，如果他们不愿意或不能够看见神圣文献中明显记载的内容，无论是圣经还是我们，都没有向他们提供他们灭亡和瞎眼的理由。因为我们都认识、阅读、相信并持守：天主是唯一的，祂创造了天地，因为我们不认识任何别的神，也永远不会认识这样的神，因为根本没有。祂说：\BibleQuote{我是天主，除我以外没有别的神，是正义的拯救者。}在另一处：\BibleQuote{我是首先的，我是末后的，除我以外没有像我的神。}以及\BibleQuote{谁曾用手心量诸水，用手虎口量苍天？谁曾用升斗盛大地，用秤称山岭，用天平称森林？}\Person{希则克雅}说：\BibleQuote{好叫普世都知道，唯独祢是天主。}再者，主亲自说：\BibleQuote{你为什么问我关于善的事？唯有天主是善的。}再者，宗徒保禄说：\BibleQuote{唯独祂是不朽的，居住在无人能接近的光中，没有人见过，也不能看见祂。}在另一处：\BibleQuote{但中保不是一方的中保，天主却是唯一的。}但是，正如我们持守、阅读并相信这一点，我们也应当不忽略圣经的任何部分，因为我们绝不应该拒绝圣经中所记载的基督神性的那些标志，免得我们因败坏圣经的权威而被视为败坏我们神圣信德的完整性。因此，让我们相信这一点，因为耶稣基督、天主之子是我们的主和天主，这是极其信实的；因为\BibleQuote{在起初已有圣言，圣言与天主同在，圣言就是天主。这圣言在起初就与天主同在。}并且\BibleQuote{圣言成了血肉，居住在我们中间。}以及\BibleQuote{我的主，我的天主。}以及\BibleQuote{他们是以色列的祖先，按血统来说，基督也是从他们而来的，祂是在万有之上，永远受赞美的天主。}那么，我们该说什么呢？圣经给我们摆出了两位天主吗？那么它怎么说\BibleQuote{天主是唯一的}？难道基督不是天主吗？那么，怎么对基督说\BibleQuote{我的主，我的天主}？因此，除非我们以适当的崇敬和合法的论证持守这一切，我们就会合理地被认为给异端者提供了绊脚石，这肯定不是由于从不欺骗的圣经的过错，而是由于人错误的妄断，他们由此选择了成为异端者。首先，我们必须转过来攻击那些指责我们说有两个天主的人。经上记载（他们不能否认）\BibleQuote{只有一位主}。那么，他们如何看待基督？——祂是主，还是根本就不是主？他们绝对不怀疑祂是主；因此，如果他们的推理正确，这里就已经有两位主了。那么，根据圣经，怎么会有\BibleQuote{只有一位主}呢？而且基督被称为\BibleQuote{唯一的师傅}。然而我们读到宗徒保禄也是一位师傅。那么，照此看来，我们的师傅就不是唯一的，因为从这些事情我们得出结论有两位师傅。那么，根据圣经，怎么会有\BibleQuote{我们的师傅只有一位，就是基督}呢？在圣经中，有一位\BibleQuote{被称为善的，就是天主}；但在同一圣经中，基督也被称为善的。那么，如果他们的结论正确，就没有一位善，而是有两位善。那么，根据圣经信德，怎么只有一位善呢？但如果他们认为，基督是主并不妨碍只有一位主的真理，保禄是我们的师傅也不妨碍唯一的师傅的真理，基督被称为善也不妨碍只有一位善的真理；那么按照同样的推理，让他们明白，从天主是唯一的事实，基督也被宣称为天主这一真理并不受阻碍。\par

% structure: p0095 source=h2 target=section reason=relative-heading-rank; base=h2

\section{第三十一章}

\Emph{但是，天主、天主之子，由永远的天主圣父所生，祂常在天父内，是圣父的第二位，祂若不奉圣父的谕旨便一无所行；而且祂是主，是天主伟大计谋的天使，圣父的神性藉本体相通而赐予祂。}\par

因此，天主圣父，万物的创立者与创造者，唯有祂不知起源，无始，不可见，无限，不死，永恒，是独一天主；对于祂的伟大、威严或权能，我不说没有任何事物能优先于祂，而是说没有任何事物能与祂相比；从祂那里，当祂愿意时，圣子、圣言诞生了；圣子并非在受击打的空气之声中，亦非在从肺腑挤出的声音音调中被领受，而是在天主所发出的权能之实体中被承认；祂神圣与天主性诞生的奥秘，宗徒未曾获悉，先知未曾发现，天使未曾知道，受造物未曾领悟。这些奥秘唯独为圣子所知晓，因为祂知道圣父的奥秘。那么，祂既由圣父所生，便常在圣父内。我这样说常，是为表明祂不是无始的，而是受生的。但那位在万时之前的，必须说祂常在圣父内；因为对于在万时之前的那位，不能指定任何时间。并且祂常在圣父内，除非圣父不常是圣父，只是圣父在某种意义上先于祂——因为必须——在某种程度上——祂在成为圣父之前必须\Emph{是}。因为那位不知起源者必须先于那位有起源者；正如祂较小，因为知道祂在祂内，因受生而有起源，并通过祂的诞生在某种程度上与圣父有相似的本性，尽管祂因受生而有起源，因为祂是由那位独无起源的圣父所生。那么，当圣父愿意时，祂从圣父发出；那在圣父内的那位，从圣父出来；那在圣父内、因为祂来自圣父的那位，随后与圣父同在，因为祂从圣父出来——就是说，那位神圣的实体，其名为圣言，万物藉祂而受造，没有祂什么也没有受造。因为万物都在祂之后，因为它们是由祂而有的。合理地，祂在万物之前，但在圣父之后，因为万物是由祂造的，而祂出自那位万物由其意志而造者。确实，天主出自天主，造成一个次于圣父的位格，作为圣子，但没有从圣父取去祂是独一天主的特性。因为如果祂不是受生的——与那无始者相比，两者都显示平等——祂就会造成两个无始的存在，从而造成两位天主。如果祂不是受生的——与那未受生者相比，并被发现平等——他们既未受生，就会合理地产生两位天主，这样基督就会成为两位天主的原因。如果祂像圣父那样无始地被形成，并且祂自身也像圣父那样是万物的开端，这就会造成两个开端，从而也向我们显示两位天主。或者，如果祂也不是圣子，而是圣父从自身生出另一位圣子，那么，合理地，与圣父相比，并被称为与祂一样伟大，祂就会造成两位圣父，从而也证明两位天主的存在。如果祂是不可见的，与那不可见者相比，并被宣告平等，祂就会展示两位不可见者，从而也证明他们是两位天主。如果是不可理解的，还有其他属于圣父的属性，合理地说，祂会引起两位天主的指控，正如这些人所虚构的。但现在，无论祂是什么，祂不是从自己，因为祂不是无始的；祂是从圣父来的，因为祂是受生的，无论是作为圣言，或是作为权能，或是作为智慧，或是作为光，或是作为圣子；无论祂是这些中的哪一个，既然祂不是从任何其他根源，如前所说，而是从圣父，把祂的起源归于祂的圣父，祂就不能因两位天主的数目在天主性中造成分歧，因为祂是从那位独一天主受生而获得祂的开端。在这种情形下，祂作为那位无始者的独生子和首生子，祂是万物的唯一根源与头。因此祂宣告天主是独一的，因为祂证明天主没有根源或开端，而更是万物的开端与根源。此外，圣子不凭自己的意愿行事，也不凭自己的决定做事；祂不是从自己来的，而是服从祂圣父的一切命令和规诫；因此，虽然出生证明祂是圣子，但顺服至死也宣告祂是祂圣父旨意的仆人，祂属于圣父。这样，祂使自己凡事顺服祂的圣父，虽然祂也是天主，却藉祂的顺服表明圣父是独一天主，祂也从圣父取得祂的开端。因此祂不能造成两位天主，因为祂没有造成两个开端，因为祂从那位无始者获得了祂在万时之前诞生的源头。因为那无始者就是其他受造物的开端——只有天主圣父是，超越那受生者的开端——而由祂所生的那位，合理地来自那无始者，证明那开端是祂自己所从出的，即使那受生者是天主，祂仍表明祂是独一天主，因为受生者证明祂是无始的。所以祂是天主，但为了这个特殊结果而受生，即祂应成为天主。祂也是主，但为了圣父的这个目的而生，即祂可成为主。祂也是天使，但祂被圣父预定为天使，以宣布天主伟大的谋略。祂的天主性如此被宣告，以免因天主性的任何不和谐或不平等而造成两位天主。因为万有由圣父作为圣子而服从于祂，而祂自己，连同那些服从于祂的，也服从于祂的圣父，祂确实被证明为祂圣父的圣子；但祂被发现为万有的主和天主。因此，当万有置于祂之下都被交付给作为天主的祂，并且万有都服从于祂时，圣子把祂所领受的一切都归于圣父，又将祂天主性的全部权柄交还给圣父。真实而永恒的圣父被显明为独一天主，唯独从祂发出这天主性的权能，也被赐予并倾注于圣子，又藉着本体的共融返回给圣父。天主确实显明为圣子，可见天主性被赐予并扩展到祂。然而，圣父仍被证明是独一天主；而在相互转移中，那威严和天主性逐步被返回并反射，如同由圣子亲自送还给那赐予它们的圣父；这样，合理地，天主圣父是万有的天主，也是祂自己所生为圣子的那位根源。此外，圣子是万有的天主，因为天主圣父将祂所生的那位置于万有之前。因此，天主与人之间的中保基督耶稣，因祂是天主，由祂自己的圣父使所有受造物的权能服从于祂；当所有受造物都服从于祂时，祂与祂的圣父天主合而为一，因祂持守那状态，即祂更\Quote{was heard,}，便简要地证明了祂的圣父是独一、唯一、真实的天主。\par

\SourceRecord{Translated by Robert Ernest Wallis.；Ante-Nicene Fathers；Vol. 5.；Edited by Alexander Roberts, James Donaldson, and A. Cleveland Coxe.；Buffalo, NY: Christian Literature Publishing Co.,；1886.；<http://www.newadvent.org/fathers/0511.htm>.}
